\chapter{Programs, Meaning, and Translation}
\label{ch:programs-meaning-translation}

A program can become faster and still become wrong.
It may return the expected number but omit a write to a file.
It may work for small inputs but fail when an integer reaches its limit.
It may pass every test we have and fail on the first input a reader tries.
The first task of a compiler writer is therefore not to make a program smaller or faster.
It is to say what the program means, and what a translation must preserve.

Consider a calculation that adds a fixed fee to an input and then multiplies the result by a fixed rate.
Suppose the input is named $x$, the fee is $7$, and the rate is written as $3+5$.
The expression is
\[
  (x+7)\times(3+5).
\]
We do not need the input to work out the rate.
Replacing $3+5$ by $8$ seems harmless.
Yet even this small change raises useful questions.
What kind of numbers are these?
Can an addition have an effect besides producing its result?
Does evaluating the expression ever fail?
Where does the replacement happen: in source text, in a tree, or in machine code?

We will answer those questions for a deliberately small language before making claims about larger ones.
The example is small enough to execute by hand, implement, test over every input, and justify with a short proof.
It gives us a working standard for the rest of the book: a transformation needs a meaning, a contract, an implementation, and evidence that the implementation keeps the contract.

\section{A small promise with a large boundary}
\label{sec:translation-promise}

A \term{compiler} translates a program from one language or representation into another.
The result need not be machine code.
A compiler may produce assembly, bytecode for a virtual machine, another source language, or an intermediate form that another tool accepts.
An \term{interpreter} executes a representation by carrying out its operations.
The same system can do both: translate one representation into another, then interpret the result.

These words describe work, not ranks of sophistication.
A tree-walking interpreter can be the clearest executable description of a language.
A large compiler may use an interpreter internally to evaluate constants.
A runtime can interpret a function while collecting the information needed to compile it later.
What matters is which representation each component accepts and which promise it makes.

For an ordinary arithmetic function, the promise might look like this:
\begin{equation}
  \operatorname{run}(C(P),i)=\operatorname{run}(P,i).
  \label{eq:deterministic-translation}
\end{equation}
Here $P$ is a source program, $C$ is a compiler, and $i$ is an allowed input.
The function $\operatorname{run}$ describes execution at the relevant level.
On the left it executes the compiled representation; on the right it executes the source representation.
Equation~\ref{eq:deterministic-translation} says that their results agree.
It does not say that their instructions, temporary values, or running times agree.

This equation has a restricted use.
It assumes that the executions being compared have results and that those results describe everything we need to preserve.
An interactive server does not simply read one input and return one number.
A program can write a file, wait forever, or receive data from another process.
We will need a richer description for such programs.
The simple equation remains useful when its assumptions are stated rather than silently extended.

The boundary also includes the environment in which compiled code runs.
If a function receives the wrong argument representation, an otherwise correct instruction sequence cannot repair the call.
If a runtime expects one object layout and the compiler emits another, every field access may be wrong.
A compiler's contract therefore depends on language rules, representation choices, and agreements with other components.
We should write those agreements down before making them difficult to change.

Appel develops compiler construction through a sequence of representations and interfaces~\cite{appel-1998-modern}.
We adopt that broad engineering lesson, but use a different starting point here: the contract of one small transformation.
The aim is to make the obligation visible before the surrounding machinery becomes large.

\section{Define the numbers before using algebra}
\label{sec:tiny-expression-domain}

Our example language uses the values
\[
  V=\{0,1,\ldots,255\}.
\]
An addition or multiplication keeps the remainder after division by $256$.
Thus $250+10$ produces $4$, not $260$.
The calculation first gives $260$ in ordinary integer arithmetic, and $260=1\times256+4$.
This is \term{modular arithmetic}, with modulus $256$.
It is a rule of the example language, not an accidental feature of the machine used to run it.

We choose a small domain for two reasons.
Every value fits comfortably on a page, and every possible input can be tested.
The wraparound also prevents us from confusing the source language's arithmetic with unbounded integer arithmetic.
A language definition should settle such a difference before an optimizer starts applying identities.

Expressions have four forms:
\begin{equation}
  e ::= \operatorname{lit}(n)
      \mid x
      \mid \operatorname{add}(e,e)
      \mid \operatorname{mul}(e,e),
  \qquad n\in V.
  \label{eq:tiny-expression-grammar}
\end{equation}
The notation after $::=$ lists the allowed shapes of an expression.
A literal contains one value $n$.
The symbol $x$ names the single runtime input.
An addition or multiplication contains two child expressions.
These forms describe a tree; they are not yet a grammar for text a user types.
There are no assignments, function calls, pointers, or loops.

We require every tree to be finite.
Every binary node has exactly two children.
Each node belongs to one place in the tree: different parents do not share a child.
These conditions may look like storage details, but the implementation will rely on them when it changes nodes in place.
They also ensure that a recursive traversal eventually reaches leaves.

Let $\llbracket e\rrbracket_a$ mean the value of expression $e$ when the input $x$ is $a$.
The brackets name evaluation; the subscript supplies the input.
The definition is
\begin{align}
  \llbracket\operatorname{lit}(n)\rrbracket_a &= n, \label{eq:tiny-literal}\\
  \llbracket x\rrbracket_a &= a, \label{eq:tiny-input}\\
  \llbracket\operatorname{add}(p,q)\rrbracket_a
    &= (\llbracket p\rrbracket_a+\llbracket q\rrbracket_a)\bmod256,
       \label{eq:tiny-add}\\
  \llbracket\operatorname{mul}(p,q)\rrbracket_a
    &= (\llbracket p\rrbracket_a\times\llbracket q\rrbracket_a)\bmod256.
       \label{eq:tiny-mul}
\end{align}
Each equation describes one tree shape.
To evaluate a binary node, evaluate its children, combine their values, and reduce the result.
Every expression has a value for every input in $V$.
There is no operation that divides by zero, dereferences an address, or calls something that may not return.

This is a \term{semantics}: a definition of what the allowed forms mean.
The grammar alone does not supply it.
We could keep the same tree shapes and choose unbounded integers, checked overflow, or floating-point values instead.
The programs would then have different meanings even though their trees looked the same.
Morgan's discussion of compiler representations emphasizes the distinction between structural rules and execution rules~\cite[Chapter~4]{morgan-1998-optimizing}.
Our arithmetic equations make the execution rules explicit for this example.

\subsection{Work out the example}
\label{sec:tiny-worked-values}

Take $a=9$.
The left addition gives $9+7=16$.
The right addition gives $3+5=8$.
The multiplication gives $16\times8=128$.
All three results already lie in $V$, so reducing them changes nothing.

Now take $a=250$.
The left addition gives a mathematical sum of $257$, whose remainder is $1$.
The right addition still gives $8$.
The multiplication gives $1\times8=8$.
The result is therefore $8$.
It would be a mistake to report the unreduced mathematical value $2056$ as the language's result.

Table~\ref{tab:tiny-example-values} records several inputs, including the wraparound boundary.
The intermediate values matter.
They show where the semantics acts, rather than merely reporting that a test passed.
For input $248$, the final multiplication is $255\times8=2040$, and $2040=7\times256+248$.
For input $249$, the left addition becomes zero, so the final result is zero.

\begin{table}
  \centering
  \caption{Each operation reduces its result modulo 256.}
  \label{tab:tiny-example-values}
  \begin{tabular}{rrrr}
    \toprule
    Input $a$ & $(a+7)\bmod256$ & $(3+5)\bmod256$ & Final value \\
    \midrule
    0   & 7   & 8 & 56 \\
    1   & 8   & 8 & 64 \\
    9   & 16  & 8 & 128 \\
    248 & 255 & 8 & 248 \\
    249 & 0   & 8 & 0 \\
    250 & 1   & 8 & 8 \\
    255 & 6   & 8 & 48 \\
    \bottomrule
  \end{tabular}
\end{table}

\section{Choose a representation that exposes the work}
\label{sec:expression-representation}

The text $3+5$ is easy for a person to recognize.
A compiler should not need to recognize it by searching for that exact string.
Spaces, parentheses, and comments can change the text without changing the calculation.
Names can also look alike without naming the same object.
Text is the right input for a parser, but usually the wrong working form for an arithmetic transformation.

An \term{abstract syntax tree} keeps the operation structure while leaving out punctuation that no longer matters.
Its nodes say which operations exist and which operands belong to them.
For our example, the root is a multiplication.
Its left child is an addition of $x$ and $7$.
Its right child is an addition of $3$ and $5$.
The tree does not need a separate node for each parenthesis.

Figure~\ref{fig:tiny-before-after} shows the original tree and the tree after replacing the constant addition.
The operation removed from runtime execution is visible as a missing internal node.
We have not changed the input-dependent addition.
We have not distributed multiplication or replaced it with shifts.
A narrowly stated change is easier to inspect than a cluster of changes performed at once.

\begin{figure}
  \centering
  \begin{tikzpicture}[level distance=12mm,
    level 1/.style={sibling distance=29mm},
    level 2/.style={sibling distance=13mm}]
    \node[basic block] at (0,0) {mul}
      child[cfg edge] {node[basic block] {add}
        child[cfg edge] {node[basic block,minimum width=8mm] {$x$}}
        child[cfg edge] {node[basic block,minimum width=8mm] {7}}}
      child[cfg edge] {node[basic block] {add}
        child[cfg edge] {node[basic block,minimum width=8mm] {3}}
        child[cfg edge] {node[basic block,minimum width=8mm] {5}}};
    \node[basic block] at (7,0) {mul}
      child[cfg edge] {node[basic block] {add}
        child[cfg edge] {node[basic block,minimum width=8mm] {$x$}}
        child[cfg edge] {node[basic block,minimum width=8mm] {7}}}
      child[cfg edge] {node[basic block] {8}};
    \node[font=\sffamily\small] at (0,0.8) {Before};
    \node[font=\sffamily\small] at (7,0.8) {After};
  \end{tikzpicture}
  \caption{Folding the right addition leaves the input-dependent subtree unchanged.}
  \label{fig:tiny-before-after}
\end{figure}

The tree is one kind of \term{intermediate representation}, or IR.
\index{IR|see{intermediate representation}}
An IR is a working language between other stages of translation.
It need not resemble the source language or the final instruction set.
Its purpose is to expose the facts a stage needs and to make the stage's assumptions checkable.
An expression tree exposes nesting.
A control-flow graph exposes possible transfers of control.
A machine-level IR exposes operations, registers, and constraints that a particular target requires.

There is no single best representation for every task.
Early in translation, preserving source structure can make diagnostics easier.
Later, explicit branches can make reachability easier to analyze.
Near code emission, explicit register uses can make storage conflicts easier to detect.
The choice of representation is therefore part of the algorithm design, not a container selected after the real work is done.

\subsection{A representation has rules}
\label{sec:expression-invariants}

A \term{representation invariant} is a condition that valid instances of a representation must satisfy.
For our trees, a literal stores a value in $V$.
An input node has no children.
An addition or multiplication has two non-null children.
No child path leads back to an ancestor.
No node has two parents.

The last two conditions are different.
A cycle makes naive recursive evaluation fail to terminate.
Sharing can terminate but still make mutation surprising.
If two parents point to one child, changing that child affects both parents.
An algorithm proved for trees cannot silently be treated as an algorithm for arbitrary graphs.

In a mature implementation, a verifier can walk an IR and check its structural rules.
Some rules are local: a binary node has two operands.
Others need wider information: an operand is defined where it is used, or a branch names an existing block.
The verifier does not prove that an optimization preserves meaning.
It checks whether the optimization has produced a legal object for the next stage to consume.

This distinction helps locate failures.
If a pass creates a malformed node, the representation contract has failed.
If the node is well formed but computes the wrong value, the semantic contract has failed.
Both are bugs, but they call for different explanations and often different tests.

\section{An interpreter as a reference point}
\label{sec:tiny-reference-interpreter}

An interpreter for this language can follow Equations~\ref{eq:tiny-literal}--\ref{eq:tiny-mul} directly.
That makes it useful as a reference implementation.
It does not make it infallible.
The interpreter can still contain a wrong operator, a missing reduction, or a bad pointer access.
Its advantage is that the correspondence between the definition and the implementation is short enough to inspect.

Listing~\ref{lst:tiny-node} gives the C node representation.
The field \code{kind} tells us how to interpret the other fields.
Only a literal uses \code{literal} as its value.
Only a binary operation uses \code{left} and \code{right} as operands.
The full header also declares the evaluator and folder and states their input contract.

\begin{listing}
  \inputminted[firstline=5,lastline=12]{c}{example-code/ch01/expr.h}
  \caption{A tagged node represents one expression form.}
  \label{lst:tiny-node}
\end{listing}

Listing~\ref{lst:tiny-evaluator} implements evaluation.
The argument \code{x} supplies the runtime input.
For a literal, the evaluator returns the stored value.
For an input node, it returns \code{x}.
For an operation, it obtains both operand values before applying the operator and reducing the result.
The assertions document important preconditions, but they are not a parser or a complete tree validator.

\begin{listing}
  \inputminted[firstline=5,lastline=20]{c}{example-code/ch01/expr.c}
  \caption{The evaluator implements the four semantic equations.}
  \label{lst:tiny-evaluator}
\end{listing}

The implementation uses C's \code{unsigned}, not an eight-bit C type.
The example language has eight-bit values; the implementation language needs room for intermediate calculations.
Each operand passed to an operation is at most $255$.
An addition is therefore at most $510$, and a multiplication is at most $65025$.
C11 guarantees that \code{unsigned int} can represent at least the values $0$ through $65535$~\cite[Section~5.2.4.2.1]{wg14-2011-c11}.
Both intermediate calculations fit before the explicit remainder operation.

This is a small but useful portability argument.
We do not depend on \code{unsigned} having exactly 32 bits, nor on a compiler providing an optional exact-width type.
We do not depend on an overflowing host calculation happening to produce the guest language's desired result.
The suffix \code{U} makes constants such as \code{256U} unsigned.
The implementation and the language definition meet at an explicit remainder operation.

The evaluator's traversal visits each node once.
If $N$ is the number of nodes and $H$ is the height of the tree, its worst-case time is $O(N)$ and its extra call-stack space is $O(H)$.
Here the height is the maximum number of nodes on a root-to-leaf path.
The bound assumes constant-time access to fields and fixed-width arithmetic.
A language with arbitrary-precision integers would need to account for the size of the numbers as well as the number of nodes.

\section{Move work from execution to translation}
\label{sec:tiny-constant-folding}

The value of $3+5$ does not depend on $x$.
Evaluating it once while preparing the program can replace repeated evaluations while running the program.
This is \term{constant folding}: evaluating a constant operation during translation and replacing it with its result.
In our language, we fold an addition or multiplication when both children have become literals.

The phrase \enquote{have become} matters.
A parent can become constant only after a child is simplified.
For example, in $(14+19)\times(6+5)$, neither operand of the multiplication is initially a literal.
After folding the additions, the multiplication has literal operands $33$ and $11$.
It can then become the literal $107$, because $33\times11=363$ and $363\bmod256=107$.

Algorithm~\ref{alg:tiny-fold} visits children before deciding whether to fold their parent.
This order is a \term{postorder traversal}.
The algorithm is written as a function that produces a replacement tree.
The C version will reuse existing node storage instead of allocating a new node for every result.
Both versions have the same semantic contract.

\begin{algorithm}
  \caption{Fold constant operations in a finite expression tree.}
  \label{alg:tiny-fold}
  \begin{pseudocode}
    \Require A finite, non-sharing tree $e$ from Equation~\ref{eq:tiny-expression-grammar}
    \Ensure A tree with the same value for every input $a\in V$
    \Function{Fold}{$e$}
      \If{$e$ is a literal or the input node}
        \State \Return $e$
      \EndIf
      \State $p\gets\Call{Fold}{\operatorname{left}(e)}$
      \State $q\gets\Call{Fold}{\operatorname{right}(e)}$
      \If{$p=\operatorname{lit}(u)$ and $q=\operatorname{lit}(v)$}
        \If{$e$ is an addition}
          \State \Return $\operatorname{lit}((u+v)\bmod256)$
        \Else
          \State \Return $\operatorname{lit}((u\times v)\bmod256)$
        \EndIf
      \EndIf
      \State \Return a node with $e$'s operator and children $p,q$
    \EndFunction
  \end{pseudocode}
\end{algorithm}

For the fee-and-rate example, the traversal first reaches $x$ and $7$.
The input is not a literal, so the left addition remains an addition.
It then reaches $3$ and $5$.
Both are literals, so the right addition becomes $8$.
The root still has a nonliteral left child, so it remains a multiplication.
The resulting expression is $(x+7)\times8$.

Originally, evaluation reached seven nodes: two addition nodes, one multiplication node, the input node, and three literals.
After folding, it reaches five: one addition, one multiplication, the input, and two literals.
The number of runtime arithmetic operations falls from three to two.
Those are statements about this interpreter and representation.
They are not yet measurements of machine-code size or elapsed time.

\subsection{One traversal is enough for this rule}
\label{sec:tiny-fold-complexity}

Once a child has been processed, it contains no operation whose two children are literals.
If the parent can become a literal, the algorithm does so immediately.
If it cannot, at least one child is nonliteral.
The parent's decision does not create an unprocessed constant operation inside that child.
There is therefore no need to repeat the traversal for this particular rule.

The algorithm does $O(1)$ local work per original node.
With $N$ nodes and height $H$, its worst-case time is $O(N)$ and its extra recursion space is $O(H)$.
The in-place implementation needs no new expression nodes.
A version that constructs a separate result tree may need $O(N)$ additional node storage.
These storage bounds should not be confused: the output representation and the traversal's temporary stack are different resources.

This argument does not say that every optimization needs only one pass.
Two different rules can expose opportunities for each other.
A change to one block can invalidate an analysis used by another.
A repeated optimization pipeline needs its own argument for termination, its own cost limit, or both.
The absence of such machinery in this example is a consequence of its narrow rule, not a general principle about compilers.

\section{Prove the replacement that was actually made}
\label{sec:tiny-fold-proof}
\index{constant folding!correctness}

Let $F(e)$ be the tree produced by Algorithm~\ref{alg:tiny-fold}.
The claim is
\begin{equation}
  \llbracket F(e)\rrbracket_a=\llbracket e\rrbracket_a
  \qquad\text{for every valid tree $e$ and every $a\in V$.}
  \label{eq:tiny-fold-correctness}
\end{equation}
Notice both quantifiers.
The claim is not limited to the fee-and-rate example, nor to the inputs in Table~\ref{tab:tiny-example-values}.
It covers every finite expression permitted by the grammar.
The small set of expression forms lets us prove it by following the tree's structure.

For a literal, $F$ returns the same literal.
Its value is unchanged by Equation~\ref{eq:tiny-literal}.
For the input node, $F$ returns the same input node.
Its value is unchanged by Equation~\ref{eq:tiny-input}.
These are the leaf cases.

Now consider an addition with children $p$ and $q$.
Assume that folding each child preserves its value for every input.
This assumption is valid because the children are smaller trees; repeatedly descending to smaller children eventually reaches the leaf cases.
Write $p'=F(p)$ and $q'=F(q)$.
For any input $a$, their values agree with those of $p$ and $q$.

If $p'$ and $q'$ are not both literals, the algorithm keeps an addition with those children.
Its value is
\[
  (\llbracket p'\rrbracket_a+\llbracket q'\rrbracket_a)\bmod256
  = (\llbracket p\rrbracket_a+\llbracket q\rrbracket_a)\bmod256.
\]
That is exactly the original addition's value.
If both children are literals, write them as $\operatorname{lit}(u)$ and $\operatorname{lit}(v)$.
The algorithm returns the literal $(u+v)\bmod256$.
The retained addition would have produced that same value by Equation~\ref{eq:tiny-add}.
Replacing it with the literal therefore preserves the value as well.

For multiplication, the argument is identical except that the local operation is multiplication and the relevant rule is Equation~\ref{eq:tiny-mul}.
Both the retained-node case and the literal-result case compute the original product modulo $256$.
The leaf cases and the two operation cases exhaust the grammar.
Equation~\ref{eq:tiny-fold-correctness} follows.
This proof method is \term{structural induction}: establish the claim for the leaves, then show how the claim for children establishes it for their parent.

\subsection{What the proof does not establish}
\label{sec:tiny-proof-boundary}

The proof is about the specified transformation on valid trees.
It does not establish that an arbitrary C program implements that transformation.
A C folder can confuse the two child pointers, write the wrong tag, or use the host machine's arithmetic incorrectly.
It can also exhaust the process stack on a sufficiently deep tree.
The mathematical function has no such process resource limit.

We therefore need a second connection: the implementation must correspond to the algorithm under stated resource and storage assumptions.
For this example, the correspondence is short enough to review directly.
Tests support that review and can expose mistakes, but they do not turn finite testing into a proof about every tree.

The distinction will remain useful as the compiler grows.
We can have a correct optimization idea with a defective implementation.
We can have a faithful implementation of an incorrect idea.
We can also have a correct pass called when its preconditions do not hold.
Each failure has a different remedy.
A proof attached to the wrong contract is not protection against the other two.

\section{Implement the fold without hiding ownership}
\label{sec:tiny-fold-implementation}
\index{constant folding!implementation}

Listing~\ref{lst:tiny-folder} is the in-place implementation.
It first leaves literals and input nodes alone.
For an operation, it folds both children.
If either child remains nonliteral, it returns without replacing the parent.
Otherwise, it evaluates the now-constant parent and changes that node into a literal.

\begin{listing}
  \inputminted[firstline=22,lastline=36]{c}{example-code/ch01/expr.c}
  \caption{The in-place folder reuses the parent node for a literal result.}
  \label{lst:tiny-folder}
\end{listing}

The call \code{eval(e, 0U)} does not claim that the actual runtime input is zero.
At that point, both children are literals.
The evaluator will read those values and will never inspect an input node.
Any input argument in $V$ gives the same result.
Using zero is merely a convenient way to satisfy the evaluator's interface.

The call is also not a hidden traversal of the original subtree.
Each child has already been folded and is now a literal.
Evaluation reaches exactly the parent and its two immediate children.
It therefore takes constant time at this point.
Calling a recursive evaluator at every node would not in general justify a linear-time bound; the local shape established by the condition is what justifies the bound here.

After obtaining the value, the code changes \code{kind} and clears both child pointers.
Clearing the pointers makes the new node's shape agree with its tag.
It does not release the old children's storage.
The caller owns all nodes, including nodes that the root can no longer reach.
The examples use stack-allocated nodes or a fixed arena, so their storage has a lifetime independent of the current tree edges.

This ownership rule must remain visible.
If each node were separately allocated and owned only through its parent, clearing the pointers without freeing the children would leak memory.
If a child were shared, freeing it could leave another parent with a dangling pointer.
Neither problem is solved by adding a call to \code{free} at the end of the function.
The representation and ownership contracts must agree with the mutation strategy.

An \term{arena} allocates a collection of objects whose storage is reclaimed together.
It is often convenient for transformations that disconnect objects during a phase.
Disconnected objects occupy space until the arena is released, so cheap reclamation is not the same as immediate reclamation.
The right choice depends on how many objects become unreachable and how long the arena remains alive.
Our fixed test arena makes that trade-off explicit without introducing a general allocator.

\section{Test the claim at several boundaries}
\label{sec:tiny-fold-testing}

The source files for this example live in \code{example-code/ch01/}.
The chapter includes code from those files rather than maintaining a second copy in LaTeX.
The directory's README gives compilation commands and the input contract.
The files use C11 and the standard library only.

The first tests check each binary operation on every pair of literal operands.
There are $256\times256=65{,}536$ pairs.
For each pair, the tests check addition and multiplication before folding, then check that folding produces the expected literal.
They also check that the replacement node no longer has children.
This is exhaustive over those two operations' literal inputs, not over arbitrary expression trees.

The second tests build the fee-and-rate tree.
They record its value for every input in $V$, fold it, and compare the results for every input again.
They check the shape as well: the root is still a multiplication, the left child is still an addition, and the right child is the literal $8$.
The shape check catches a folder that does nothing but happens to preserve all results.
Semantic preservation alone does not establish that an optimizer performed its intended work.

The tests then fold the result again and compare the values again.
An operation is \term{idempotent} when applying it twice has the same result as applying it once.
For this rule, a second traversal has no additional constant operation to fold.
Repeated application should therefore leave the reachable tree unchanged.
The current test checks repeated evaluation results; a stronger test could also compare the complete reachable shape.
We should distinguish the mathematical property from the part of it a particular test checks.

The third group builds 256 deterministic generated trees of height at most five nodes on any root-to-leaf path.
For each tree, it compares all 256 input results before and after folding.
The generator uses a fixed arithmetic rule, so a failure can be reproduced.
These cases provide additional shapes, but they do not cover all trees of that height or prove correctness for all heights.
The test set is finite even when the property is not.

The program \code{demo.c} prints selected before-and-after values.
Its output agrees with the corresponding rows of Table~\ref{tab:tiny-example-values}.
The tests also check the nested constant expression used in the folding discussion, the wraparound sum $250+10$, and the two identities used in the exercises.
Hand calculations, representation checks, and executable comparisons support one another; none should be allowed to stand in for all the others.

\subsection{The reference can share a mistake}
\label{sec:tiny-oracle-risk}

Before-and-after evaluation can miss a mistake shared by both executions.
If the evaluator accidentally subtracts where it should add, the folder may use that same evaluator to compute a constant.
The two versions can then agree with each other while both disagree with the language definition.
Agreement is useful evidence only when we also ask what the compared implementations have in common.

This is why the direct literal-operation tests compute expected values separately from the evaluator's traversal.
The arithmetic expectations still use C, so this is not a fully independent implementation of arithmetic.
The earlier range argument is part of the justification for those expectations.
The hand-worked boundary cases add another check on the intended meaning.

For larger languages, a second compiler, a separately written interpreter, or a translation validator can provide a different point of comparison.
Each has limits.
Two compilers can share assumptions, and a validator can have its own bugs.
The useful question is not whether a component deserves the label \enquote{trusted} in isolation.
It is which failures the evidence can detect and which components remain outside that evidence.
Safonov discusses compiler trustworthiness as a broader engineering concern than the correctness of one optimization~\cite{safonov-2010-trustworthy}.

\section{A result is not the whole execution}
\label{sec:observable-behavior}

Our expression language is unusually forgiving.
Every expression terminates and has no effect except producing a value.
That makes result equality a sufficient comparison.
Adding one operation with an effect changes the problem.

Suppose a function increments a counter and returns $7$.
Multiplying its return value by zero produces zero, but calling the function also changes the counter.
Listing~\ref{lst:zero-effect} is a complete C program that makes the difference visible.
The function call is evaluated before the multiplication has its result.
The final print observes both the value and the updated counter.

\begin{listing}
  \inputminted{c}{example-code/ch01/observations.c}
  \caption{A zero result does not make the function call's effect dispensable.}
  \label{lst:zero-effect}
\end{listing}

The program prints \code{value=0 calls=1}.
Replacing the entire body of \code{keep\_effect} with a return of zero would leave the counter unchanged.
The program would print a different result and return a failure status.
A compiler may remove the multiplication while retaining the increment.
It may also inline the function or reorganize the surrounding calculation.
It must preserve the program's required observations, not its particular spelling or call structure.

An \term{observable behavior} is an aspect of execution that the language and its execution model require an implementation to preserve.
The exact definition belongs to the language.
It is not safe to invent a universal list from an optimizer's current needs.
For C, the abstract-machine rules discuss observable requirements including volatile accesses, file contents at termination, and interactive input/output behavior~\cite[Section~5.1.2.3]{wg14-2011-c11}.
Our counter matters because its value contributes to printed output; an arbitrary unobserved local update need not be preserved as a machine operation.

The word \enquote{observable} does not mean that a debugger can inspect a temporary while single-stepping.
Debugging interfaces can impose additional expectations, but those expectations are a separate agreement.
Nor does it mean every physical property a measuring device could detect.
Ordinary language semantics generally allows changes in instruction timing; a security-sensitive system may need stronger rules about timing or memory-access patterns.
The preservation boundary must match the promise being made.

\subsection{Record a trace when order matters}
\label{sec:behavior-traces}

A \term{trace} records a sequence of events.
For a small input/output language, events might be \enquote{read this value} and \enquote{write this value}.
Two runs that write the same values in different orders have different traces.
Two runs that both return zero can also have different traces if only one writes a message.

We can describe a terminating deterministic execution by a pair $(t,r)$, where $t$ is its trace and $r$ is its final result.
A more complete model can distinguish a return, a defined trap, and divergence.
The details depend on which constructs the language contains.
An expression language with no loops needs less machinery than a concurrent language with shared memory and external inputs.

The point of introducing a trace is not to make every proof longer.
It is to prevent the proof from forgetting behavior the program can actually have.
If a transformation only touches a pure expression inside a larger program, we may prove a local value property and then justify why that local replacement preserves the surrounding execution.
If the expression can perform input/output, the local property must include that behavior or explicitly require its absence.

\section{Equality, permission, and refinement}
\label{sec:behavior-refinement}

Some language rules allow more than one behavior for the same program and input.
For example, a language might leave the evaluation order of two independent operations open.
A compiler can choose one permitted order without supporting every order on every execution.
Requiring the compiled program to reproduce the entire set of source possibilities would then be unnecessarily strong.

Let $\mathcal{B}(P,i)$ be the set of behaviors permitted for program $P$ on input $i$.
For programs whose behavior is defined, a useful compiler requirement is
\begin{equation}
  \mathcal{B}(C(P),i)\subseteq\mathcal{B}(P,i).
  \label{eq:behavior-refinement}
\end{equation}
Here both sides must describe observations at a common interface.
Machine-level bookkeeping events should not appear on only one side of the comparison.
Equation~\ref{eq:behavior-refinement} says that a compiled execution must not introduce a forbidden behavior.
This form of restriction is called \term{refinement}.

The equation alone is not a complete definition of compiler correctness.
If an execution model allowed a target program to have no behaviors at all, an empty set would satisfy the inclusion vacuously.
We also need the model and correctness statement to account for execution, progress, and termination in the intended way.
For a terminating deterministic source program, describing divergence as a distinct behavior prevents an endlessly looping target from masquerading as a correct result.

Undefined behavior needs particular care.
A language can decline to impose execution requirements after a program violates a specified rule.
That is not the same as choosing one of several ordinary defined results.
It is also not a guarantee that the program will stop with an error message.
The scope of the compiler's obligation must be stated on inputs and executions to which the language gives a meaning.

In C, signed integer overflow is not a request for wraparound arithmetic.
The exceptional-expression rule makes behavior undefined when a result is not representable in its type, subject to the language's specific rules for operations and types~\cite[Section~6.5, paragraph~5]{wg14-2011-c11}.
Unsigned arithmetic has a different rule: results are reduced modulo one more than the largest representable value of the resulting unsigned type~\cite[Section~6.2.5, paragraph~9]{wg14-2011-c11}.
Integer promotions still affect which type an expression uses.
Naming an eight-bit storage type does not by itself make every intermediate C expression an eight-bit operation.

Our example avoids those complications by specifying its own domain and proving that the host operations fit.
The lesson is not that every compiler should use this arithmetic.
It is that a transformation must follow the arithmetic of the language or IR it is transforming.
An identity copied from a mathematics textbook is not a substitute for that check.

\section{Control flow changes what can happen}
\label{sec:first-control-flow-graph}

An expression tree tells us how operands nest.
It does not directly express a choice between paths through a program.
Suppose a program tests whether the input is below $128$.
On the true path it computes $(x+7)\times8$; on the false path it returns $x$.
Only one path executes on each run.

A \term{control-flow graph}, or CFG, represents possible transfers of control.
\index{CFG|see{control-flow graph}}
Its vertices are blocks of instructions, and its directed edges describe possible transfers between them.
A \term{basic block} is a sequence of instructions entered at its beginning, with no internal branch that selects a different successor before its end.
For this first graph, the instructions are just descriptive labels.
We will define an executable instruction representation before implementing graph transformations.

Figure~\ref{fig:first-diamond} shows the choice as a diamond.
The entry block tests the input.
Its two outgoing edges have different conditions.
Each path assigns the result to $r$, then reaches a return block.
The graph describes both possible paths even though an individual execution follows only one of them.

\begin{figure}
  \centering
  \begin{tikzpicture}[node distance=14mm and 12mm]
    \node[basic block] (test) {test $x<128$};
    \node[basic block,below left=of test] (yes) {$r\gets(x+7)\times8$};
    \node[basic block,below right=of test] (no) {$r\gets x$};
    \node[basic block,below right=of yes] (ret) {return $r$};
    \draw[cfg edge] (test) -- node[above left,font=\small] {true} (yes);
    \draw[cfg edge] (test) -- node[above right,font=\small] {false} (no);
    \draw[cfg edge] (yes) -- (ret);
    \draw[cfg edge] (no) -- (ret);
  \end{tikzpicture}
  \caption{A branch selects one computation; the return block is shared.}
  \label{fig:first-diamond}
\end{figure}

For $x=9$, the true edge is selected and the result is $128$.
For $x=250$, the false edge is selected and the result is $250$.
The fee-and-rate expression would evaluate to $8$ at that second input, but this execution does not evaluate it.
Distinguishing an expression's possible value from whether it executes is essential to reasoning about a whole program.

If we moved an operation from one branch to the entry block, it could run on inputs that previously skipped it.
For our finite, pure arithmetic expression, that extra evaluation would terminate and would have no external effect.
For a memory access, a division, or a function call, the extra evaluation could matter.
Moving work requires a justification about execution conditions as well as about the operation's returned value.

The CFG itself does not tell us every semantic fact we need.
An edge says a transfer is possible in the representation.
An analysis may prove that a particular edge cannot be taken under a particular condition.
Until such a proof exists, an optimizer cannot simply omit a path because the path looks unlikely.
Likelihood supports a performance decision; impossibility supports a different kind of correctness argument.

\section{Separate facts from rewrites}
\label{sec:analysis-and-transformation}

A \term{program analysis} computes facts used to reason about a program.
A \term{program transformation} changes the program's representation.
Our folder does both in a small local form: it recognizes two literal children and replaces their parent with a literal.
Larger optimizations often separate those tasks so that several rewrites can reuse the same analysis.

Suppose a value is known to be $8$ at one instruction.
An analysis can record that fact.
A transformation can replace a multiplication's operand with the literal $8$.
Another transformation can simplify a comparison using the same fact.
The analysis is useful because its statement has a defined scope: which value, which program point, and which possible executions.

Using a fact creates three distinct obligations.
The analysis must cover the executions that can reach the point, the rewrite must be valid under the reported fact, and the implementation must apply that rewrite to the representation the fact describes.
For example, knowing that \code{x} is zero does not by itself justify deleting the evaluation of \code{f() * x}: the call may have an observable effect or may fail to return.
Even a sound fact and a valid rule are insufficient if an intervening edit has changed which definition of \code{x} reaches the multiplication.
This separation gives a practical way to diagnose wrong code: inspect the fact, the rule's preconditions, and the fact's lifetime before changing the optimization heuristic.

The cost of a wrong fact can be much larger than the cost of a missed opportunity.
If an analysis cannot establish that a value is constant, keeping the original computation may make the program slower.
If it falsely reports a constant, replacing the computation can make the program wrong.
A \term{conservative analysis} deliberately keeps enough possibilities to avoid making an unsound claim.
The direction of that conservatism depends on what the analysis represents.
For possible targets of a call, retaining an extra target can be safe while omitting a real target is not.

The small folder makes no claim about the value of $x$.
It is possible that every caller supplies the same input, but our function's contract allows all $256$ values.
Without an additional argument about callers, the folder must retain the input-dependent work.
Not exploiting an unstated assumption is a form of correctness, not a weakness to conceal.

\subsection{A change can invalidate a fact}
\label{sec:analysis-invalidation-introduction}

An analysis result refers to a particular program representation.
After a transformation changes that representation, the old result may no longer apply.
If a graph gains a new predecessor edge, a fact justified by all previous predecessors may become false.
If a store is inserted, an earlier conclusion about a loaded value may stop being valid.

There are three broad responses.
We can recompute the analysis.
We can update it while changing the program.
Or we can discard it and avoid using it until it is computed again.
Each choice has an implementation cost.
Continuing to use the old result without justification is not a fourth choice.

This is one reason to give passes explicit interfaces.
A pass should state which representation it expects, which properties it preserves, and which analyses it invalidates.
Those statements let a pipeline compose passes without depending on accidental knowledge inside each implementation.
The interface is part of the correctness argument even when the source code expresses it only through assertions and documentation.

\section{A compiler is a chain of small languages}
\label{sec:compiler-phase-contracts}

Figure~\ref{fig:first-pipeline} sketches a translation path.
The diagram is a map of responsibilities, not a rule that every compiler must have exactly four stages.
The tree preserves source structure.
The middle representation exposes operations and control.
The machine representation exposes target constraints.
Each arrow crosses a boundary where one set of invariants must be established from another.

\begin{figure}
  \centering
  \begin{tikzpicture}[node distance=8mm]
    \node[basic block,minimum width=24mm] (source) {Source};
    \node[basic block,minimum width=24mm,right=of source] (tree) {Typed tree};
    \node[basic block,minimum width=24mm,right=of tree] (ir) {Middle IR};
    \node[basic block,minimum width=28mm,right=of ir] (machine) {Machine IR};
    \draw[cfg edge] (source) -- (tree);
    \draw[cfg edge] (tree) -- (ir);
    \draw[cfg edge] (ir) -- (machine);
  \end{tikzpicture}
  \caption{Different representations make different obligations explicit.}
  \label{fig:first-pipeline}
\end{figure}

The \term{front end} checks and translates the source language.
It typically handles tokens, syntax, names, and types, then creates a representation suitable for later work.
The \term{middle end} performs analyses and transformations at a level that need not commit to every detail of a target instruction set.
The \term{back end} maps work onto a target's operations, storage resources, and calling conventions.
These names describe a useful division of responsibility, not three sealed rooms.

A source feature can create obligations all the way to machine code.
If an operation must report overflow, lowering cannot silently replace it with an operation that merely wraps.
If a call can trigger garbage collection, the runtime may need a way to find live references at that call.
If a function can unwind, instruction motion must respect the defined behavior of unwinding.
The representation can change, but an obligation cannot disappear simply because its original syntax has disappeared.

A \term{lowering} translates a construct into operations that express it more explicitly or are closer to the eventual execution mechanism.
For example, a structured conditional can become a test and branch edges.
A representation with a single allocation operation may later become a call plus runtime metadata.
Lowering is not inherently an optimization.
It can increase the number of visible operations while making the required implementation possible.

At each boundary, we should ask two questions.
Does the result satisfy the next representation's structural rules?
Does it preserve the source construct's required behavior?
The first question is often suitable for an IR verifier.
The second needs reasoning about the translation itself.
Keeping the questions separate makes a complicated pipeline easier to test and explain.

\section{Optimization needs an objective, not a name}
\label{sec:optimization-objectives}

An \term{optimization} aims to improve a program according to some objective while preserving the required behavior.
The word does not mean that the compiler finds the best possible program.
Many passes make local choices or apply heuristics under limited information.
The objective might be execution time, code size, memory use, energy, startup delay, or some combination.
Two objectives can disagree about the same change.

Our folder removes one runtime addition from the example expression.
If the same expression is interpreted many times, avoiding that repeated work is a plausible benefit.
If the expression is interpreted once, preparing the folded tree also costs work.
Whether the change reduces total time depends on how preparation and execution are used.
We should not turn an operation count into a benchmark result.

For a simple model, let $T$ be the extra translation time spent optimizing a program.
Let $s$ be the execution time saved per run, and let $k$ be the number of runs.
Ignoring other differences, the optimization pays for its extra preparation time when
\begin{equation}
  k s > T.
  \label{eq:optimization-break-even}
\end{equation}
The units must agree: if $T$ and $s$ are times, $k$ is a count.
This is a model for a decision, not a measured property of the folder.
If $s$ is zero or negative, increasing the run count does not produce the assumed benefit.

An ahead-of-time compiler and a runtime compiler can face different values of $k$ and different tolerances for $T$.
A build system may spend substantial time preparing a binary that runs many times.
A runtime may need to decide whether a particular function is worth compiling while the application is already running.
The same transformation can be appropriate under one budget and inappropriate under another.

Code size introduces another trade-off.
Duplicating a block can avoid a branch on one path while making the program larger.
Inlining can expose constants while increasing the amount of code later passes must process.
The claim \enquote{this pass improves the program} is incomplete until the objective and relevant conditions are named.
Morgan's treatment of optimizing compilers connects transformations to machine and compiler-engineering concerns rather than treating them as isolated algebra~\cite[Chapters~1--2]{morgan-1998-optimizing}.

\section{Make failures small enough to understand}
\label{sec:compiler-failure-method}

When a compiler produces a wrong answer, the original failing program can be much larger than the cause.
A useful first goal is a smaller program that still exposes the error.
The smaller program makes the disputed contract easier to identify.
It can also become a regression test: an executable reminder of a failure the compiler must not repeat.

For our folder, a wrong result at input $250$ suggests several possibilities.
The evaluator may have omitted a reduction.
The folder may have calculated a constant with the wrong operator.
The test may have built the tree incorrectly.
Printing the input, the original tree, the transformed tree, and both values can separate those possibilities without printing every internal pointer.

The relevant observation should be part of the test.
For a pure expression, the result value is enough.
For Listing~\ref{lst:zero-effect}, the test must also check the counter-dependent output or final status.
A test that records only the returned arithmetic value cannot detect a lost effect by inspecting that value more often.

The allowed inputs also belong in the test description.
A supposed compiler bug in an arbitrary C program may instead depend on undefined behavior in that program.
That does not make every surprising output correct; it means the language rules must be checked before the example can establish a miscompilation.
Using a deliberately defined small language lets us practice this discipline without first confronting every detail of C.

During development, structural checks close to a transformation can catch damage before it spreads.
For this example, assertions on node kinds and literal ranges help document assumptions.
A complete graph implementation will need broader checks for edges and definitions.
When the first invalid state is found close to the pass that created it, the final machine-code failure becomes less mysterious.

\section{Translate the tree into a sequence}
\label{sec:tiny-stack-translation}

The constant folder changes a tree into another tree.
We can check its result with the same evaluator because the representation has not changed.
A compiler must eventually cross a harder boundary: the result and the input may have different shapes and different execution rules.
We need a way to explain why those different executions still agree.

Our first such translation removes the tree from execution.
It produces a sequence of instructions for a small stack machine.
The machine does not follow child pointers.
It reads instructions in order and keeps intermediate values on a stack.
The input still belongs to $V$, and every arithmetic result still belongs to $V$.
We change the execution mechanism without changing the arithmetic promise.

A \term{stack machine} obtains operands from a stack rather than naming a separate storage location for every intermediate value.
The top is the end from which values are removed.
For our notation, the top appears on the right.
Thus $[23,7,8]$ has $8$ at its top.
The sequence $[23,7]$ remains after removing that value.
This convention is small enough to state once and easy enough to get wrong if left unstated.

The instruction set has four operations.
\code{PUSH} carries a literal operand and pushes its value.
\code{GET\_INPUT} pushes the current input.
\code{ADD8} removes two values, adds them modulo $256$, and pushes the result.
\code{MUL8} does the same with multiplication.
There are no branch instructions and no explicit return instruction.
Reaching the end with exactly one stack value completes a program.
That value is its result.

The compiler's task is now concrete.
For every valid expression, it must produce instructions that finish with its value.
It must not consume a value that the expression's caller already had on the stack.
It must not rely on an instruction that reads a nonexistent operand.
These are related obligations, but they are not interchangeable.

\subsection{Give each instruction a meaning}
\label{sec:tiny-stack-semantics}

Write $S\mathbin{\|}[v]$ for the stack obtained by appending $v$ to stack $S$.
The stack $S$ can be empty.
For input $a$, the instruction rules are
\begin{align*}
  \code{PUSH}\ n &: S \longmapsto S\mathbin{\|}[n],\\
  \code{GET\_INPUT} &: S \longmapsto S\mathbin{\|}[a],\\
  \code{ADD8} &: S\mathbin{\|}[u,v]
    \longmapsto S\mathbin{\|}[(u+v)\bmod256],\\
  \code{MUL8} &: S\mathbin{\|}[u,v]
    \longmapsto S\mathbin{\|}[(u\times v)\bmod256].
\end{align*}
The arithmetic instructions need at least two values.
Trying to execute one with fewer values is \term{stack underflow}.
We reject such instruction sequences before running them.
This is an error in the target representation, not a newly introduced result of the source expression language.

For a binary instruction, $v$ is removed first because it is on top.
The value $u$ is removed second.
The calculation uses $u$ as its left operand and $v$ as its right operand.
Addition and multiplication give the same answer if those two operands are exchanged.
A later subtraction instruction would not.
We keep the operand convention explicit even though the present instruction set cannot expose that particular mistake through a changed numeric answer.

Let $\operatorname{Exec}(c,a,S)$ mean executing instruction sequence $c$ with input $a$ and initial stack $S$.
For now this notation describes the mathematical machine, with enough storage and a valid instruction sequence.
Execution of an empty sequence leaves $S$ unchanged.
Execution of a nonempty sequence applies its first instruction and then executes the remaining sequence.
The finite sequence and absence of jumps guarantee that this execution reaches its end.
The bounded C implementation will report failure when it cannot meet its storage contract.

\subsection{Generate instructions in operand order}
\label{sec:tiny-stack-generation}

Write $C(e)$ for the instruction sequence compiled from expression $e$.
The translation rules are
\begin{align*}
  C(\operatorname{lit}(n)) &= [\code{PUSH}\ n],\\
  C(x) &= [\code{GET\_INPUT}],\\
  C(\operatorname{add}(p,q)) &= C(p)\mathbin{\|}C(q)\mathbin{\|}[\code{ADD8}],\\
  C(\operatorname{mul}(p,q)) &= C(p)\mathbin{\|}C(q)\mathbin{\|}[\code{MUL8}].
\end{align*}
Here the same concatenation symbol joins instruction sequences rather than stacks.
The surrounding expressions identify which kind of sequence is involved.
These rules choose left-to-right operand evaluation.
That order is permitted because the source expressions are pure and total.
It is a choice made by this compiler, not a general permission to reorder effectful operations in other languages.

Algorithm~\ref{alg:tiny-stack-emit} gives the corresponding traversal.
Emission appends one instruction to a growable logical sequence.
The mathematical algorithm assumes enough storage; the implementation checks a fixed capacity.
Like folding, emission visits children before completing their parent.
Unlike folding, it does not replace the parent with another tree node.
It appends an instruction that will combine the values left by the children's code.

\begin{algorithm}
  \caption{Emit stack instructions for an expression.}
  \label{alg:tiny-stack-emit}
  \begin{pseudocode}
    \Require A valid expression tree $e$; an instruction sequence with enough capacity
    \Ensure Append $C(e)$ without changing the existing instruction prefix
    \Function{Emit}{$e$}
      \If{$e=\operatorname{lit}(n)$}
        \State append $\code{PUSH}\ n$
      \ElsIf{$e=x$}
        \State append $\code{GET\_INPUT}$
      \Else
        \State \Call{Emit}{$\operatorname{left}(e)$}
        \State \Call{Emit}{$\operatorname{right}(e)$}
        \State append $\code{ADD8}$ or $\code{MUL8}$ according to $e$'s operator
      \EndIf
    \EndFunction
  \end{pseudocode}
\end{algorithm}

Our running expression produces seven instructions.
The input and $7$ must be pushed before their addition.
The literals $3$ and $5$ must then be pushed before the second addition.
Finally, the multiplication combines the two subtree results.
Table~\ref{tab:tiny-stack-trace} works through the complete sequence at input $250$.
The first addition wraps to $1$, just as it did in the tree evaluator.

\begin{table}
  \centering
  \caption{The seven-instruction program produces 8 at input 250.}
  \label{tab:tiny-stack-trace}
  \begin{tabular}{rll}
    \toprule
    Position & Instruction & Stack after execution \\
    \midrule
    & Before execution & $[\,]$ \\
    0 & \code{GET\_INPUT} & $[250]$ \\
    1 & \code{PUSH} 7 & $[250,7]$ \\
    2 & \code{ADD8} & $[1]$ \\
    3 & \code{PUSH} 3 & $[1,3]$ \\
    4 & \code{PUSH} 5 & $[1,3,5]$ \\
    5 & \code{ADD8} & $[1,8]$ \\
    6 & \code{MUL8} & $[8]$ \\
    \bottomrule
  \end{tabular}
\end{table}

Figure~\ref{fig:tiny-stack-boundary} separates the representations involved.
The source evaluator and target executor reach the same value through different data structures.
The compiler does not run the target program to discover its answer.
It builds an instruction sequence that will work for every allowed input.
At input $9$, the same seven instructions produce $128$.
There is no input-specific substitution in the generated program.

\begin{figure}
  \centering
  \begin{tikzpicture}[node distance=15mm and 17mm]
    \node[basic block] (tree) {Expression tree};
    \node[basic block,right=of tree] (code) {Stack code};
    \node[basic block,below=of tree] (value1) {Value in $V$};
    \node[basic block,below=of code] (value2) {Value in $V$};
    \draw[cfg edge] (tree) -- node[above,font=\small] {compile} (code);
    \draw[cfg edge] (tree) -- node[left,font=\small] {evaluate} (value1);
    \draw[cfg edge] (code) -- node[right,font=\small] {execute} (value2);
    \draw (value1) -- node[above,font=\small] {equal} (value2);
  \end{tikzpicture}
  \caption{Correct translation connects two execution mechanisms at a common result.}
  \label{fig:tiny-stack-boundary}
\end{figure}

\section{Prove preservation of the surrounding stack}
\label{sec:tiny-stack-proof}
\index{stack machine!correctness}

Starting from an empty stack is how we run a whole expression program.
It is not enough for the most useful induction statement.
When the right subtree starts, the left subtree's value is already on the stack.
The right subtree's code must preserve that value while producing its own.
A proof restricted to an empty stack cannot directly justify that step.

Instead, we prove the stronger statement
\begin{equation}
  \operatorname{Exec}(C(e),a,S)
  = S\mathbin{\|}[\llbracket e\rrbracket_a]
  \label{eq:tiny-stack-prefix}
\end{equation}
for every valid expression $e$, every input $a\in V$, and every finite initial stack $S$ of values in $V$.
The execution must also avoid underflow.
Equation~\ref{eq:tiny-stack-prefix} says that compiled expression code preserves the entire initial stack and appends one result.
It states exactly the property that a parent needs from each child.

For a literal, the emitted instruction pushes that literal and changes nothing below it.
For the input node, the emitted instruction pushes $a$ and changes nothing below it.
These cases establish the statement for leaves.
Neither instruction removes an operand, so neither can underflow.

For an addition, assume the statement holds for children $p$ and $q$.
Executing $C(p)$ on $S$ gives $S\mathbin{\|}[u]$, where $u=\llbracket p\rrbracket_a$.
Now use the claim for $q$ with that larger initial stack.
Executing $C(q)$ gives $S\mathbin{\|}[u,v]$, where $v=\llbracket q\rrbracket_a$.
Both child executions avoid underflow by the induction assumptions.
The final \code{ADD8} has the two operands it needs.
It replaces them with $(u+v)\bmod256$ without touching $S$.
By the source addition rule, this value is $\llbracket\operatorname{add}(p,q)\rrbracket_a$.

For multiplication, the two child executions establish the same stack shape.
The final \code{MUL8} replaces $u,v$ with their product modulo $256$.
That is the source multiplication's value.
All expression forms have now been covered, so structural induction establishes Equation~\ref{eq:tiny-stack-prefix}.
Choosing $S$ empty gives the whole-program correctness statement.

The stronger statement was not extra decoration.
It let the proof follow the implementation's actual composition rule.
When an induction stalls because a recursive call runs in a larger context, the missing ingredient can be a claim about preserving that context.
Here the context is an existing stack prefix.
In a different compiler proof, it might be untouched memory, an environment of names, or a relation between source and target stores.
The appropriate claim depends on what the generated code can read or change.

\subsection{Do not confuse a fragment with a program}
\label{sec:tiny-stack-fragments}

Compiled expression code has a net stack effect of plus one.
That is true whether it starts with zero values or twenty.
A whole expression program additionally starts empty and finishes with exactly one value.
Those extra conditions belong to its execution interface.

For example, take a stack containing $23$ and execute the running expression's code at input $250$.
The result is $[23,8]$.
Nothing has gone wrong: the fragment preserved its caller's value and appended its result.
But if we present that two-value final stack as a complete expression program's result, we have violated the whole-program interface.

The C compiler can append a fragment to a code buffer that already contains instructions.
It therefore cannot promise that every successful append produces a complete executable program.
The existing prefix might leave values on the stack or might itself be malformed.
The compiler promises how it appends code; the verifier checks the resulting whole sequence.
Separating these promises prevents an append operation from being mistaken for an endorsement of arbitrary existing code.

\section{Check the code before executing it}
\label{sec:tiny-stack-verifier}
\index{stack machine!verification}

A target instruction sequence can have a plausible length and still be invalid.
An arithmetic instruction can appear before its operands.
A literal can lie outside $V$.
An instruction tag can be unknown.
The sequence can finish with zero values or leave several values instead of one.
We can detect all of these faults without knowing the runtime input.

The verifier scans the instructions while tracking only stack depth.
A push increases depth by one.
An arithmetic instruction requires depth at least two and decreases it by one.
At the end, depth must be exactly one.
The verifier also records the largest depth it encounters.
This is an exact bound for this branch-free machine: every accepted instruction executes on every run and its stack effect is independent of input values.

We also choose a canonical representation for unused instruction operands.
Only \code{PUSH} uses its argument field.
The argument field of every other opcode must contain zero.
This rule is not necessary for the arithmetic meaning, but it avoids several encodings of the same instruction caused by irrelevant leftover data.
It gives tests and diagnostic tools a simpler representation to inspect.
The compiler establishes the rule; the verifier checks it rather than trusting the compiler to have done so.

The verifier does not prove that the program implements a particular source tree.
The sequence consisting only of \code{PUSH} with operand $99$ is valid.
It is still the wrong translation of most source expressions.
Structural safety and semantic preservation remain different properties even when both are easy to test.

\subsection{Bound the storage explicitly}
\label{sec:tiny-stack-storage}

The implementation in \code{stack.h} reserves room for $64$ instructions.
This number is an engineering limit of the example, not part of the mathematical source language.
The code structure contains an instruction array and a logical length named \code{used}.
The valid portion is the prefix ending just before \code{used}.
The remaining slots are storage, not instructions in the program.

Because every push consumes an instruction, an accepted program of at most $64$ instructions cannot need more than $64$ stack slots.
The executor therefore uses a fixed stack array of that size.
The verifier rejects a logical code length larger than its instruction array before reading any instruction.
It checks every instruction before execution starts.
The executor then uses the verified stack effects to avoid an underflow or out-of-bounds stack write.

This argument assumes that the code object is a real, initialized C object and that its active instructions do not change between verification and execution.
It does not turn arbitrary bytes received from an untrusted file into a valid C structure.
There is no byte decoder in this example.
There is also no concurrent writer changing the instruction array during the run.
A serialized format or concurrent runtime would need additional boundary rules.

The compiler takes a limit no greater than the array capacity.
If an append would exceed the limit, it reports failure.
The public wrapper restores the original logical length when any recursive emission fails.
The existing active instruction prefix remains unchanged.
Unused slots may contain partially emitted instructions, but they lie outside the restored logical length.
This is a limited rollback guarantee, not a promise to restore every byte of the code object.

For the original seven-node example, an empty buffer with limit seven succeeds.
Limit six fails, and the buffer's logical length returns to zero.
After constant folding, the five-node expression needs five instruction slots.
Its code sequence pushes the input, pushes $7$, adds, pushes $8$, and multiplies.
The arithmetic result remains unchanged even though the capacity needed by this target representation falls.

\section{Read the compiler and executor together}
\label{sec:tiny-stack-c-implementation}

Listing~\ref{lst:tiny-stack-emitter} shows the recursive emission helper from \code{stack.c}.
The two recursive calls are separate statements.
Their left-to-right order is therefore visible in the implementation rather than depending on the evaluation order of function arguments in C.
The helper stops as soon as an append fails.
The public wrapper, shown in the source file, handles restoration of the logical length.

\begin{listing}
  \inputminted[firstline=13,lastline=25]{c}{example-code/ch01/stack.c}
  \caption{Emission follows the source tree while appending target instructions.}
  \label{lst:tiny-stack-emitter}
\end{listing}

Listing~\ref{lst:tiny-stack-runner} gives the executor.
It validates the input range and the whole code sequence before running any instruction.
The variable \code{sp} is a count of occupied stack slots.
Consequently, the next pushed value goes into \code{stack[sp]}, and removing a value first decrements \code{sp}.
Keeping that convention consistent is more useful than relying on a variable name such as \enquote{top}, which could mean either an index or a count.

\begin{listing}
  \inputminted[firstline=66,lastline=85]{c}{example-code/ch01/stack.c}
  \caption{The executor runs verified code and publishes one result on success.}
  \label{lst:tiny-stack-runner}
\end{listing}

The executor removes the right value into \code{b} and then the left value into \code{a}.
Its unsigned intermediate arithmetic has the same range argument as the tree evaluator.
Both operands lie in $V$, so the largest product is $65025$.
The explicit remainder implements the target instruction's specified arithmetic.
The implementation does not rely on a host overflow to stand in for the guest machine's reduction.

The output parameter is assigned only after successful verification and execution.
On rejection, the caller's existing result value remains unchanged.
Failure is reported through a Boolean status rather than a special numeric result.
Every value in $V$, including zero, is a legitimate program result.
Using zero to mean both a valid answer and a failed execution would make the interface needlessly ambiguous.

\subsection{Account for both time and space}
\label{sec:tiny-stack-costs}

For a tree with $N$ nodes and height $H$, successful compilation emits exactly $N$ instructions.
Each leaf emits one push and each binary node emits one arithmetic instruction.
With a preallocated instruction buffer and constant-time append, compilation takes $O(N)$ time and $O(H)$ additional recursion space.
The target code occupies $O(N)$ slots.
This statement describes the algorithm as capacity grows; the example's fixed limit restricts which instances it accepts.

The verifier takes $O(N)$ time and $O(1)$ auxiliary space when instruction count is $N$.
The executor takes another $O(N)$ time.
This implementation verifies on every call, so a reported execution cost includes both scans.
A runtime could verify an immutable code object once and execute it many times, but it would then need a trustworthy association between the checked object and the object being executed.
Skipping the scan without establishing that association would remove the premise of the stack-safety argument.

The exact peak stack demand can be calculated from the source tree.
Let $D(e)$ be the maximum number of extra stack slots needed while executing $C(e)$.
For either leaf, $D(e)=1$.
For either binary operation with children $p,q$,
\begin{equation}
  D(e)=\max\bigl(D(p),\,1+D(q)\bigr).
  \label{eq:tiny-stack-demand}
\end{equation}
The first term covers evaluation of the left child.
The second covers evaluation of the right child while the left result occupies one slot.
The final arithmetic instruction does not increase the peak.

For the original running expression, each addition has demand two.
The root therefore has demand $\max(2,1+2)=3$.
After folding the right addition into a literal, the demand becomes $\max(2,1+1)=2$.
The verifier obtains the same peaks from the instruction sequences without seeing the original trees.
Agreement between the two calculations is another useful implementation check.

\section{Compose folding with translation}
\label{sec:tiny-stack-composition}

We now have two distinct preservation results.
Folding preserves the source expression's value.
Compilation preserves that value while moving from tree evaluation to stack execution.
The two results can be connected because folding produces a valid tree of the kind compilation expects.
That representation condition is essential: equality of values would not excuse a malformed tree passed to the compiler.

With enough target capacity, the combined claim is
\begin{align*}
  \operatorname{Exec}(C(F(e)),a,[\,])
    &= [\llbracket F(e)\rrbracket_a]\\
    &= [\llbracket e\rrbracket_a].
\end{align*}
The first equality comes from the compiler's stack-prefix theorem.
The second comes from the folder's value-preservation theorem.
Neither proof needs to know the internal implementation of the other pass.
Their contracts identify the shared boundary.

The bounded implementation adds a status distinction to this account.
A translation that reports insufficient capacity has not produced an executable replacement.
The preservation claim concerns successful translations under the stated resource assumptions.
The caller must handle failure rather than execute unused or partially written array slots.
In a larger compiler, that handling might grow a buffer, select another representation, or report a controlled compilation failure.
It must not quietly reinterpret a resource failure as the source program's answer.

\subsection{Test the connected claim}
\label{sec:tiny-stack-testing}

The new test program, \code{test\_stack.c}, compares source evaluation with target execution before and after folding.
For each selected tree, it records all $256$ source results before mutation.
It compiles the original tree and compares the target result at every input.
It then folds the tree, compiles again, and compares the second target program against the saved source results.
Saving the original answers prevents the test from accidentally treating the already folded tree as its only reference.

The running example gets additional structural checks.
The test checks all seven original opcodes and their operands, not merely the final numeric result.
It checks peak depth three before folding and peak depth two afterward.
It checks that the folded program has five instructions.
Those assertions expose a translation that preserves answers but fails to implement the intended representation or simplification.

For the arithmetic instructions, the tests cover both operations on all $65{,}536$ literal operand pairs.
For translation across larger shapes, they generate $256$ bounded trees and check every input before and after folding.
Those generated cases remain a finite sample of source shapes.
Their value is in testing the connections among components, not in replacing the structural proof with a large count of successful runs.

Rejection tests exercise the target interface separately.
They include an empty program, an arithmetic instruction with no operands, an out-of-range literal, an unknown opcode, nonzero unused operands, excess final values, and an impossible logical length.
They also check that a failed run leaves its output parameter alone.
Capacity tests exercise partial emission and restoration of the logical prefix.
A nonempty-prefix test confirms that a successful fragment append need not constitute a valid whole program.

These tests cover a boundary that result comparisons alone would miss.
A compiler used only on valid small trees might never emit an invalid instruction.
Its verifier still needs rejection tests because the verifier's contract is broader than the compiler's normal output set.
Conversely, a verifier that accepts a program says nothing about whether the compiler chose the correct program for its source.
Testing the components together does not remove the need to test their separate promises.

\section{Make a wrong translation fail on purpose}
\label{sec:tiny-fault-lab}

A passing test suite tells us that the tested programs met the tested conditions.
It does not tell us whether those conditions would notice the mistake we fear.
We can ask a more concrete question: if multiplication were translated as addition, which test would fail?
This is not a proposal to damage the working compiler.
It is a controlled experiment with a separate, deliberately faulty component.

The program \code{fault\_lab.c} uses the correct compiler to emit stack code, then changes every \code{MUL8} instruction to \code{ADD8}.
It leaves each instruction's argument field alone.
It also leaves the source tree alone.
The fault lives only in this laboratory program; \code{compile\_expr} in \code{stack.c} remains unchanged.
Keeping the specimen separate makes it harder to confuse a demonstration of a bug with the implementation the reader should use.

Listing~\ref{lst:tiny-mutant} shows the faulty wrapper.
Both arithmetic instructions require two stack values and produce one.
Both require a zero argument field.
Replacing one with the other therefore preserves the structural properties checked by our verifier.
The generated program remains executable by the defined stack machine, but it generally computes the wrong value.

\begin{listing}
  \inputminted[firstline=7,lastline=15]{c}{example-code/ch01/fault_lab.c}
  \caption{This isolated wrapper deliberately changes multiplication into addition.}
  \label{lst:tiny-mutant}
\end{listing}

The specimen illustrates a narrow form of \term{mutation testing}: introduce a controlled change and check whether the tests detect its consequence.
One mutation does not model every compiler failure.
It checks a particular weakness in the evidence.
Here the question is whether we compare arithmetic meaning, rather than merely accepting structurally legal code.
The distinction fits the broader separation between compiler assurance tasks discussed by Safonov~\cite{safonov-2010-trustworthy}.
The specimen and its reduction procedure here are original, small experiments, not a survey of that subject.

\subsection{Choose a witness with a known answer}
\label{sec:tiny-fault-witness}

Use this source expression:
\[
  ((x+7)\times(3+5))+11.
\]
Its tree has nine nodes.
It extends the earlier seven-node tree with a final addition and the literal $11$.
At input $9$, the correct calculation is
\[
  ((9+7)\times(3+5))+11=16\times8+11=139.
\]
Every intermediate result fits in $V$, so no reduction changes these numbers.
The faulty target instead adds $16$ and $8$, then adds $11$, producing $35$.
We have a concrete disagreement: the expected value is $139$, and the target value is $35$.

A \term{counterexample} to a universal preservation claim is one allowed source program and input for which that claim fails.
The word does not require a large random test or an exotic boundary value.
Our ordinary input $9$ is enough.
The source is valid, its operations are total, and the target passes structural verification.
None of those facts can remove the disagreement.

It is useful to freeze the input while investigating this failure.
If a reducer changes the program and the input together, it may keep finding failures without showing which change simplified the explanation.
Our first reduction procedure keeps the witness input at $9$.
Later, once the program is smaller, we can ask whether the input matters at all.
These are two separate simplification decisions.

\section{Define which failure a reducer must retain}
\label{sec:tiny-reduction-predicate}

A \term{test-case reducer} searches for a smaller input that still exposes a chosen failure.
It needs a precise predicate for that failure.
For this laboratory, let $M(e,a)$ hold when all of the following are true:
\begin{enumerate}
  \item The faulty compiler successfully emits the whole program for $e$.
  \item The emitted code passes structural verification.
  \item The executor successfully runs it at input $a$.
  \item The returned value differs from $\llbracket e\rrbracket_a$.
\end{enumerate}
We test source validity by construction: every candidate is made from valid tree pieces or valid literal leaves.
The tree remains finite and non-sharing.
The predicate therefore describes wrong-code behavior on a valid, executable case.
It does not count a compilation failure or a rejected target as the same failure.

Listing~\ref{lst:tiny-mismatch} implements the predicate.
Returning false after a compilation failure does not claim that the compiler is correct.
It says that this candidate does not satisfy the particular predicate we are preserving.
Different investigations need different predicates.
An investigation of a compiler crash should retain that crash, not quietly substitute a wrong numeric result for it.

\begin{listing}
  \inputminted[firstline=17,lastline=27]{c}{example-code/ch01/fault_lab.c}
  \caption{The predicate retains successful executions with disagreeing values.}
  \label{lst:tiny-mismatch}
\end{listing}

This distinction prevents a common reduction mistake.
Suppose deleting part of a source program introduces an undeclared name.
If any nonzero tool status counts as interesting, the reducer may preserve a syntax error instead of the original miscompilation.
The final file can then be smaller but useless for explaining the bug.
Our structured candidates avoid syntax errors, and our predicate rejects changes that no longer reach a wrong-code execution.

The predicate also relies on a reference.
Here that reference is the previously defined evaluator, with arithmetic checked independently on literal pairs and small worked examples.
If the evaluator were wrong, a disagreement would need further investigation rather than an automatic declaration that the compiler was at fault.
Reduction makes the evidence easier to inspect; it does not certify the reference implementation.

\section{Reduce a tree without breaking its ownership}
\label{sec:tiny-greedy-reducer}

The reducer considers four kinds of replacement at each node.
It can replace a binary node with its left child, replace it with its right child, or replace any node with the literal $0$ or $1$.
It tries these choices in that order.
After an accepted replacement, it restarts its search at the root.
If no replacement at the current node is accepted, it searches its children from left to right.

This is a \term{greedy algorithm}: it keeps the first acceptable local improvement rather than examining every possible smaller tree before deciding.
It is not an algorithm for finding the smallest possible counterexample in an unrestricted language.
The candidate choices and their order matter.
They define which smaller cases this search can reach.

We also need a measure that decreases whenever a replacement is accepted.
Node count alone is not sufficient because replacing the input leaf with a literal keeps the node count unchanged.
Let $N(e)$ be the number of reachable tree nodes.
Let $W(e)$ be the sum of the weights of its leaves, assigning weight $n$ to literal $n$ and weight $256$ to the input leaf.
Compare the pair $(N(e),W(e))$ lexicographically: prefer fewer nodes, then a smaller weight sum when the node counts agree.

Every accepted replacement must decrease this pair and retain $M(e,9)$.
Replacing a binary node by a child decreases node count.
Replacing an input by $0$ or $1$ decreases weight without changing node count.
Replacing literal $3$ by literal $1$ decreases weight as well.
Replacing literal $0$ by literal $1$ is not a decrease and is not even tested for preservation.
Both components are nonnegative integers, so a strictly decreasing sequence of these pairs cannot continue forever.

Algorithm~\ref{alg:tiny-greedy-reduce} describes the outer search.
The candidate enumeration restarts after each successful replacement because the tree and its useful sites have changed.
The algorithm terminates when none of its permitted single replacements both reduces the measure and retains the predicate.
That stopping condition is narrower than global minimality.

\begin{algorithm}
  \caption{Reduce a failing tree using a fixed set of local replacements.}
  \label{alg:tiny-greedy-reduce}
  \begin{pseudocode}
    \Require A valid finite tree $e$ with $M(e,a)$; fixed input $a$
    \Ensure Retain $M(e,a)$; no permitted single improving replacement remains
    \Repeat
      \State $\mathit{changed}\gets\mathit{false}$
      \ForAll{sites in preorder, trying child replacements then literals $0,1$}
        \State form a candidate tree $e'$
        \If{$(N(e'),W(e'))<(N(e),W(e))$ and $M(e',a)$}
          \State $e\gets e'$
          \State $\mathit{changed}\gets\mathit{true}$
          \State \Break
        \EndIf
      \EndFor
    \Until{not $\mathit{changed}$}
  \end{pseudocode}
\end{algorithm}

The implementation changes one node's stored structure temporarily.
It saves the original structure, installs the proposed replacement, and evaluates the predicate on the entire root.
If the candidate fails, it restores the saved structure.
Predicate evaluation does not fold or otherwise mutate the tree.
Consequently, restoring the one changed structure is enough to restore the previous reachable tree.

Replacing a node with a child's structure follows the caller-owned storage rule established earlier.
The child's storage remains alive.
Disconnected nodes are not freed during reduction.
The original laboratory tree uses stack storage that outlives every predicate call, and no candidate introduces a cycle or a shared reachable node.
This is not a general-purpose reducer for arbitrary graph ownership schemes.
Its storage reasoning depends on precisely the same non-sharing contract as the earlier folder.

For the nine-node case, the C counters used by the measure are comfortably within the minimum range of unsigned integers.
An unbounded reduction library would need a separate argument or checked arithmetic for those counters.
The executable specimen is a bounded laboratory, not an unchecked promise about trees of arbitrary size.

\section{Work through every accepted reduction}
\label{sec:tiny-reduction-walkthrough}

The first candidate replaces the root addition with its left child.
The trailing $11$ disappears.
The new expression is the earlier product, which still disagrees at input $9$: the correct value is $128$, and the faulty value is $24$.
This step reduces node count from nine to seven.

At the product root, neither child alone retains the failure.
Both children contain only addition and leaves, so the faulty wrapper changes no instruction in their translations.
Replacing the root with $0$ or $1$ also removes the fault.
Those candidates are rejected.
The search moves into the left subtree instead.

Replacing $x+7$ with its left child gives $x\times(3+5)$.
At the fixed input, the correct value is $9\times8=72$ and the faulty value is $9+8=17$.
The candidate is accepted.
Replacing the input leaf with literal $0$ then gives $0\times(3+5)$.
The correct result is $0$; the faulty result is $8$.
Node count stays five in this second step, but leaf weight decreases.

The right subtree can next be replaced by its left child, literal $3$.
The result $0\times3$ computes $0$ in the source and $3$ in the faulty target.
Replacing $3$ with $0$ would erase the disagreement: both $0\times0$ and $0+0$ produce zero.
That candidate is rejected.
Replacing $3$ with $1$ retains the disagreement and decreases weight, so it is accepted.

Table~\ref{tab:tiny-reduction-steps} records the full accepted path.
The emitted diagnostic program prints the same node counts, input, and value pairs.
We can therefore check the hand reasoning against an executable trace of the search rather than merely asserting that reduction succeeded.

\begin{table}
  \centering
  \caption{Five accepted replacements retain the wrong-code failure at input 9.}
  \label{tab:tiny-reduction-steps}
  \begin{tabular}{lrrr}
    \toprule
    Source expression & Nodes & Correct & Faulty \\
    \midrule
    $((x+7)\times(3+5))+11$ & 9 & 139 & 35 \\
    $(x+7)\times(3+5)$ & 7 & 128 & 24 \\
    $x\times(3+5)$ & 5 & 72 & 17 \\
    $0\times(3+5)$ & 5 & 0 & 8 \\
    $0\times3$ & 3 & 0 & 3 \\
    $0\times1$ & 3 & 0 & 1 \\
    \bottomrule
  \end{tabular}
\end{table}

The final expression no longer contains the input.
It disagrees for every allowed input, not merely for the original witness.
Its correct translation is a push of zero, a push of one, and multiplication.
The faulty translation ends with addition instead.
Both sequences pass the verifier, both reach a stack of depth two before the final instruction, and both finish with one value.
The only relevant difference is the meaning of that final operation.

Figure~\ref{fig:tiny-fault-diagnosis} places the two conclusions side by side.
A structural acceptance is correct: there is no malformed instruction or missing operand.
A semantic acceptance would be wrong: zero and one are distinct results in $V$.
Adding more checks of the same stack-depth property cannot discover that difference.
The evidence must cross the source-to-target semantic boundary.

\begin{figure}
  \centering
  \begin{tikzpicture}[node distance=13mm and 17mm]
    \node[basic block] (source) {$0\times1$};
    \node[basic block,below left=of source] (good) {PUSH 0; PUSH 1\\MUL8};
    \node[basic block,below right=of source] (bad) {PUSH 0; PUSH 1\\ADD8};
    \node[basic block,below=of good] (zero) {valid; result 0};
    \node[basic block,below=of bad] (one) {valid; result 1};
    \draw[cfg edge] (source) -- node[above left,font=\small] {correct} (good);
    \draw[cfg edge] (source) -- node[above right,font=\small] {faulty} (bad);
    \draw[cfg edge] (good) -- (zero);
    \draw[cfg edge] (bad) -- (one);
  \end{tikzpicture}
  \caption{The verifier accepts both programs; only one preserves the source value.}
  \label{fig:tiny-fault-diagnosis}
\end{figure}

\subsection{Say exactly how small the result is}
\label{sec:tiny-reducer-minimality}

For this particular mutant, a literal or an input node cannot fail.
Its compiled code contains no multiplication instruction for the wrapper to change.
A source tree with an internal binary node has at least three nodes.
The reduced counterexample has exactly three nodes.
It is therefore smallest by node count for this mutant within our source grammar.

That conclusion uses a separate argument about the fault.
The greedy algorithm itself did not establish it simply by stopping.
For another fault, a different local replacement or a coordinated change at two sites could reveal a smaller failing program.
Local irreducibility means that the allowed single-step candidates do not improve the current case.
Global minimality means that no smaller failing case exists in the entire stated search space.
The words describe different claims and require different evidence.

Even a smallest node count need not give a unique explanation.
Several three-node multiplications expose this mutation.
Our tie-breaking rule prefers small literal weights, and the fixed candidate order leads to $0\times1$.
A reducer that tries literal $2$ first could stop at a different representative if its rules did not later reduce that literal.
The exact output belongs to the specified search procedure, not to an unspecified notion of the one true minimal test.

\section{Repair the contract boundary, then keep the test}
\label{sec:tiny-fault-regression}

The reduced program points directly to the broken local step.
The compiler's preservation proof requires a source multiplication to end in a target multiplication.
Replacing it with addition changes the value equation used in that induction case.
The stack-prefix shape remains correct, but its appended value is no longer the promised value.
The failure is in the arithmetic correspondence, not in the proof's storage reasoning.

In the working compiler there is nothing to repair: the injected fault is confined to the laboratory wrapper.
The program checks that the correct compiler runs the reduced case successfully and agrees with source evaluation for all inputs.
It also checks that the mutant disagrees for all inputs.
Both assertions matter to the demonstration.
Without the first, the laboratory could merely compare two defective paths.
Without the second, the specimen might stop demonstrating the intended fault after an unrelated edit.

A \term{regression test} preserves a previously exposed failure as an executable check against its return.
The small expression is a good candidate because its expected answer can be justified without the original larger tree.
It should complement, not replace, the earlier literal-pair tests.
Those tests already cover multiplication of zero by one and many nearby operand pairs.
The laboratory explains why that coverage detects a semantic error which the verifier correctly accepts as structurally safe.

Keep the pipeline configuration with the test.
The faulty wrapper changes emitted code, and the laboratory deliberately compiles the original source tree without invoking the folder.
If we first folded $0\times1$ into literal zero, no \code{MUL8} would be emitted.
The faulty wrapper would have nothing to change, and the disagreement would disappear.
That disappearance would not show that the faulty translation was repaired.
It would show that the selected execution path no longer reaches the broken operation.

This is why a useful failure record includes the source, input, enabled passes, target configuration, expected observation, and actual observation.
It also includes whether compilation and structural verification succeeded.
For this case all arithmetic and configuration choices fit in a small standalone C program.
In a larger compiler, the same discipline prevents a minimized file from losing the flags or pass ordering needed to reproduce its original failure.

\subsection{Know what reduction costs}
\label{sec:tiny-reduction-cost}

A predicate evaluation recompiles and executes a candidate and also evaluates its source tree.
With $N$ source nodes and enough capacity, each of those steps takes $O(N)$ time here.
Computing a candidate's size measure is another tree traversal.
The number of candidate attempts depends on the tree, the accepted steps, and the search order.
It would be misleading to call the whole reducer linear merely because one predicate call is linear.

If there are $R$ complete search sweeps over trees of at most $N$ nodes, each sweep considers at most a constant number of replacements at each node.
Using the simple full-tree predicate, this gives an $O(RN^2)$ upper bound for those sweeps, with fixed-width arithmetic and constant-capacity append treated as constant-time operations.
The bound does not state that every candidate takes the worst-case time.
It states which repeated work the straightforward implementation can perform.
The decreasing measure proves termination, but does not promise a small value of $R$ for an arbitrary starting tree.

For our nine-node example, the exact accepted path is more informative than an asymptotic slogan.
There are five accepted replacements and a final unsuccessful sweep.
We print only accepted states so that the demonstration stays readable.
A debugging tool could additionally log rejected candidates and their reason for rejection.
Those records should distinguish a nondecreasing measure, a lost mismatch, and a failure to compile or execute.
They are different explanations for why a proposed reduction was not retained.

The completed experiment gives us a useful division of labor.
The verifier guards target structure.
The source evaluator and target executor expose a semantic disagreement.
The reducer makes that disagreement small enough to inspect.
The proof identifies the local correspondence that was broken.
The regression test keeps the discovered case available after the investigation ends.
No one of these components supplies all the evidence that the others provide.

\section{Total operations end where traps begin}
\label{sec:trap-total-boundary}

Every expression in Equation~\ref{eq:tiny-expression-grammar} terminates with a value in $V$.
It has no other outcome.
That property is what makes result equality a sufficient comparison for the folder and the stack compiler.
An operation with that property is \term{total}: it is defined for every combination of its operand values.
An operation that also has no observable effect is \term{pure}.
Our tiny language is pure and total, and both preservation proofs depend on those two facts.
This section introduces division with a defined trap and separates its failure behavior from externally visible effects.

A \term{trap} is a defined failure that stops ordinary value production and reports that no value was produced.
It is not a value in $V$.
It is not undefined behavior in the C sense.
It is a behavior the language definition names, such as division by zero, an out-of-range access, or an explicit assertion failure.
The implementation must detect the condition and take the defined failure path rather than continuing with an arbitrary number.
\index{trap}
\index{total operation}
\index{partial operation}

Consider an extended expression form $\operatorname{div}(p,q)$.
Its intended meaning is integer division of the dividend by the divisor.
When the divisor is zero, there is no quotient to return.
We therefore give the form two possible outcomes for each input.
Let $\llbracket e\rrbracket_a$ be either a value in $V$ or the token $\text{trap}$.
Literals and input nodes still return values.
Addition and multiplication now propagate a child trap; only when both children return values do their original arithmetic equations apply.
For the new form, first evaluate both children left to right.
If either child traps, the division traps without further calculation.
If both succeed with values $u$ and $v$, then the division traps when $v=0$ and otherwise returns the mathematical quotient $u/v$ truncated toward zero.
\begin{equation}
  \llbracket\operatorname{div}(p,q)\rrbracket_a=
  \begin{cases}
    \text{trap} & \text{if $\llbracket p\rrbracket_a=\text{trap}$},\\
    \text{trap} & \text{if $\llbracket q\rrbracket_a=\text{trap}$},\\
    \text{trap} & \text{if $q$ evaluates to $0$},\\
    u/v & \text{otherwise, with $u,v$ the two values}.
  \end{cases}
  \label{eq:trap-div-semantics}
\end{equation}
The order of the first two cases matters once effects are present, so we state it now even though either order gives the same trap-or-value answer for this pure extension.
% TODO(cite): verify the C division and trap wording against N1570 clauses before citing a paragraph number.
The C rules for division are checked against the committee draft rather than assumed from algebra~\cite{wg14-2011-c11}.

The quotient needs no modulo reduction in this domain.
Both operands lie in $V$, so the largest dividend is $255$ and the smallest nonzero divisor is $1$.
Every quotient therefore lies in $V$ already.
The wraparound still matters for computing the operands.
The divisor $249+7$ evaluates to $(249+7)\bmod256=0$, and a division by that divisor traps even though neither literal is zero on its own.
A reader who checks only the source literals will miss the trap.

Table~\ref{tab:trap-div-cases} works through the cases the laboratory checks exhaustively.
The first two rows fold only pure additions before dividing.
The third row folds a divisor to zero and must preserve the trap.
The fourth row shows truncation rather than wraparound.
The fifth row is the ordinary fee-and-rate shape with a safe divisor.

\begin{table}
  \centering
  \caption{Division traps exactly when its evaluated divisor is zero.}
  \label{tab:trap-div-cases}
  \begin{tabular}{llll}
    \toprule
    Dividend & Divisor & Folded divisor & Outcome \\
    \midrule
    $x+7$ at $x=9$ ($16$) & $3+5$ ($8$) & $8$ & $2$ \\
    $x+7$ at $x=250$ ($1$) & $3+5$ ($8$) & $8$ & $0$ \\
    $x+7$ & $249+7$ & $0$ & $\text{trap}$ for every input \\
    $7$ & $2$ & $2$ & $3$ \\
    $0$ & $5$ & $5$ & $0$ \\
    \bottomrule
  \end{tabular}
\end{table}

The second row deserves attention because its quotient is zero for a different reason than a trap.
At input $250$, the dividend is $1$ and the divisor is $8$, so integer division gives $0$.
That zero is a value: execution continues normally.
A test that treats every zero as a failure would confuse this case with the trapping row above it.
The laboratory therefore distinguishes the Boolean success flag from the numeric result, and rejection tests check that a failed run leaves its output parameter unchanged.

Listing~\ref{lst:trap-evaluator} implements the extended evaluator for this laboratory specimen.
It does not change \code{expr.c} or \code{stack.c}.
It follows Equation~\ref{eq:trap-div-semantics} directly, including left-to-right propagation of a child trap.
The return value reports success or trap, and the output parameter carries a value only on success.
This separation is the same discipline used by the stack executor, where every value in $V$ is a legitimate answer and failure needs its own channel.

\begin{listing}
  \inputminted[firstline=19,lastline=48]{c}{example-code/ch01/trap_lab.c}
  \caption{The laboratory evaluator propagates traps left to right.}
  \label{lst:trap-evaluator}
\end{listing}

The host arithmetic argument for this file mirrors the earlier one for addition and multiplication, with one added guard.
Each operand passed to division lies in $V$, so the mathematical quotient is at most $255$.
C11 guarantees that \code{unsigned int} covers at least $0$ through $65535$~\cite[Section~5.2.4.2.1]{wg14-2011-c11}.
The quotient therefore fits before any conversion.
The implementation uses \code{unsigned} rather than a narrow type, so the operands do not undergo a surprising promotion to a signed type.
Most importantly, the code tests \code{b == 0U} before evaluating \code{a / b}.
Unsigned division by zero has no defined result in C, so executing it to \enquote{see what happens} would not be a measurement but an appeal to undefined behavior.
The guard makes the laboratory's trap a defined branch on a checked condition, not an execution of an undefined operation.
The exhaustive test checks all $256\times256=65{,}536$ literal pairs and confirms that exactly the $256$ pairs with divisor zero report a trap.

Folding must now consult the divisor before replacing anything.
Folding $3+5$ to $8$ inside $(x+7)/(3+5)$ is still safe, because the divisor is a nonzero constant and the division itself still depends on $x$.
Folding the whole division to a literal would be wrong, because its dividend is not constant.
Folding $(249+7)$ to $0$ is safe as a substep, but the resulting division $(x+7)/0$ must remain a division.
Replacing it with any literal in $V$ would turn a program that traps on every input into a program that returns normally.
Listing~\ref{lst:trap-guarded-folder} shows the guard: a division whose folded children are both literals folds only when the divisor literal is nonzero.
Otherwise the node is left in place for a later stage to handle as a trapping operation.

\begin{listing}
  \inputminted[firstline=50,lastline=70]{c}{example-code/ch01/trap_lab.c}
  \caption{A division by a folded zero divisor is left unfolded.}
  \label{lst:trap-guarded-folder}
\end{listing}

The laboratory checks both sides of that guard on every input.
For the always-trapping tree $(x+7)/(249+7)$, source evaluation traps for all $256$ inputs before folding.
After folding, the tree has the shape $\operatorname{div}(x+7,0)$ and still traps for all $256$ inputs.
For the safe tree $(x+7)/(3+5)$, folding produces divisor $8$ and evaluation succeeds for all $256$ inputs with value $((a+7)\bmod256)/8$.
At $a=9$ this gives $16/8=2$, and at $a=250$ it gives $1/8=0$.
Hand calculation, exhaustive literal-pair checks, and whole-tree comparisons support one another here, as they did for the pure folder.
None of them alone establishes the guarded rule.

The stack machine from Section~\ref{sec:tiny-stack-translation} cannot express this new behavior without an extension.
Its four instructions always produce a value when they satisfy the verifier.
It has no failure status for a trapping arithmetic operation and no way to report \enquote{this program correctly stopped}.
Compiling a trapping source expression to that target would therefore require one of two explicit decisions.
Either the compiler rejects programs that may trap, with a diagnostic that names the untranslatable construct, or the target is extended with a checked instruction and an execution interface that reports failure separately from values.
What the compiler must not do is emit an unchecked host division and hope the hardware trap coincides with the source language's defined trap.
That hope confuses two different contracts.

\section{Effects make order part of meaning}
\label{sec:trap-effects-order}

A pure operation can be evaluated in any order without changing the program's observations.
An \term{effect} removes that freedom.
Writing to a file, updating a counter that later output depends on, or sending a message all create observations beyond the returned number.
Listing~\ref{lst:zero-effect} already showed that multiplying a call's result by zero does not make the call dispensable.
This section makes the ordering consequence concrete with a small event log.
\index{effect}
\index{observable behavior!effects}

The laboratory log records one tag per call in the order the calls execute.
Listing~\ref{lst:trap-log} shows the helper: it appends \code{tag} to the log and returns \code{value} unchanged.
The recorded sequence is the trace discussed in Section~\ref{sec:behavior-traces}.
Two executions can return the same number with different traces, and a program that prints or checks the log can observe the difference.

\begin{listing}
  \inputminted[firstline=72,lastline=83]{c}{example-code/ch01/trap_lab.c}
  \caption{Logging preserves the value while recording the call order.}
  \label{lst:trap-log}
\end{listing}

Consider two calls that return $10$ and $20$.
Evaluating the first call and then the second records $[1,2]$ and returns $30$.
Evaluating the second call and then the first records $[2,1]$ and still returns $30$.
The laboratory checks exactly this: both orders give $30$, but the tag sequences differ in the predicted positions.
A transformation that swaps the two calls preserves the sum while changing the trace.
Whether that change is allowed depends on what the language promises about the log.
If the log models file output, the swap reorders bytes in the file.
If the log models a debugging counter that no later observation reads, a later chapter's analysis might justify the swap, but the justification must name the unobserved property explicitly.

This is why the stack compiler's left-to-right choice was flagged as a choice rather than a general permission.
For the pure total language, evaluating $p$ before $q$ or $q$ before $p$ gives the same value and no trace, so the compiler may pick either order and document it.
For an effectful language, the source definition must state the order or state that both orders are permitted refinements in the sense of Equation~\ref{eq:behavior-refinement}.
The compiler may then choose an order only where the definition allows it.
Copying the left-to-right emission template into an effectful compiler without that check is a common way to inherit a subtle bug.

Duplication and deletion follow the same principle.
Unrolling duplicates effectful instructions in the representation, but must preserve the original dynamic trace, including iteration order and the number of effects.
Inlining a function that writes to a file moves that write to the call site, which is correct only if the language permits the write to occur there.
Dead-code elimination must not delete a call whose effect contributes to an observation, even when its returned value is unused.
The pure folder's deletion of $3+5$ was safe because the deleted operations had no trace by definition.
The zero-multiplication identity in Exercise~\ref{ex:tiny-zero-identity} is safe in the tiny language for the same reason and unsafe for the C program in Listing~\ref{lst:zero-effect} for the opposite reason.
The identity as arithmetic does not decide the matter; the effect contract does.

A useful discipline is to keep the value reasoning and the trace reasoning in separate but connected statements.
The value claim says what number is produced when execution succeeds.
The trace claim says which events occur and in which order.
The guarded division above needed only the first kind of extension, because its failure mode produces no trace beyond the trap itself.
An output operation needs the second kind, because reordering two successful executions can still be wrong.
Later analyses will exploit this separation by tracking values and effects with different abstractions rather than forcing one lattice to describe both.

\section{Nontermination is a behavior, not a missing value}
\label{sec:trap-divergence}

Every program in the tiny language finishes.
Its trees are finite, its evaluator recurses on smaller subtrees, and its stack machine executes a finite instruction sequence with no jumps.
Real languages admit executions that never finish.
A loop may never meet its exit condition, a recursive function may never reach its base case, or two threads may wait for each other forever.
Such an execution is said to \term{diverge}.
Divergence is a behavior of the program, not an absent entry in a table of results.
A compiler that turns a diverging program into a terminating one, or vice versa, has changed what the program does even if no wrong number ever appears.
\index{divergence}
\index{termination}

The laboratory cannot demonstrate divergence by actually diverging, since a test that never finishes proves nothing about the next transformation.
Instead it models bounded work with an explicit fuel parameter.
Listing~\ref{lst:trap-fuel} runs a loop that needs exactly \code{needed} steps and reports whether \code{fuel} was sufficient.
With \code{needed}=5 and \code{fuel}=3 it returns failure and leaves its output untouched.
With \code{needed}=5 and \code{fuel}=5 it succeeds and reports $5$.
With \code{needed}=0 and \code{fuel}=0 it succeeds immediately.
The tests check all three boundaries.

\begin{listing}
  \inputminted[firstline=85,lastline=97]{c}{example-code/ch01/trap_lab.c}
  \caption{Fuel turns potential divergence into a checked failure.}
  \label{lst:trap-fuel}
\end{listing}

Fuel is a testing device, not a semantics of the source language.
It lets us check the surrounding logic---preservation of success, failure, and output discipline---without waiting forever.
It does not prove that a particular source loop diverges, nor does it assign a value to a divergent execution.
Its value is narrower: it forces the implementation to separate \enquote{finished with this answer} from \enquote{did not finish within the allowance}, in the same way the division laboratory separates a zero value from a trap.
A common bug in both settings is to reuse zero for both outcomes and then test only the number.

Divergence constrains motion in the same way traps and effects do.
Figure~\ref{fig:trap-hoist} shows the pattern once, since the same diagram covers all three.
On the left, a branch guards an operation that may trap, write, or diverge.
On inputs that take the other edge, the operation never executes.
On the right, the operation has been hoisted above the test so that it executes on every input.
If the operation is pure and total, the hoisted program still returns the same values, and the change may be a reasonable performance decision under the cost model of Section~\ref{sec:optimization-objectives}.
If the operation may trap on the previously skipped path, the hoisted program introduces a failure the source never had.
If it writes, the hoisted program extends every trace with an event the source omitted on that path.
If it diverges on that path, the hoisted program fails to reach code the source reached.
Each variant breaks refinement in Equation~\ref{eq:behavior-refinement} by adding a behavior the source set did not contain.

\begin{figure}
  \centering
  \begin{tikzpicture}[node distance=14mm and 12mm]
    \node[basic block] (test) {test $c$};
    \node[basic block,below left=of test] (yes) {risky $op$};
    \node[basic block,below right=of test] (no) {skip};
    \node[basic block,below right=of yes] (ret) {return $r$};
    \draw[cfg edge] (test) -- node[above left,font=\small] {true} (yes);
    \draw[cfg edge] (test) -- node[above right,font=\small] {false} (no);
    \draw[cfg edge] (yes) -- (ret);
    \draw[cfg edge] (no) -- (ret);
    \node[basic block,right=38mm of test] (op2) {risky $op$};
    \node[basic block,below=of op2] (test2) {test $c$};
    \node[basic block,below=of test2] (ret2) {return $r$};
    \draw[cfg edge] (op2) -- (test2);
    \draw[cfg edge] (test2) -- (ret2);
    \node[font=\sffamily\small] at (1.2,1.0) {Before};
    \node[font=\sffamily\small] at (6.6,1.0) {After};
  \end{tikzpicture}
  \caption{Hoisting a risky operation above its guard adds behaviors the source lacked.}
  \label{fig:trap-hoist}
\end{figure}

The converse motion has the opposite risk.
Sinking an operation into one branch removes its effect from the other path.
If the operation diverges in the source on inputs that now skip it, the transformed program terminates where the source did not.
If it traps there, the transformed program returns where the source reported failure.
If it writes there, the transformed program omits an event the source produced.
Whether sinking is allowed therefore depends on a positive fact about the removed behavior on that path: that it cannot trap, cannot write, and must terminate, or that the language permits its removal.
An analysis that establishes such a fact is useful precisely because it discharges one of these obligations.
Without the fact, the motion is merely hopeful.

Our tiny language sidesteps all of this by construction, which is why it is a good place to practice the vocabulary.
Every operation terminates, produces a value, and leaves no trace.
Hoisting, sinking, duplication, and deletion change only the number of interpreter steps, not the observable pair of trace and result.
The guarded division, the event log, and the fuel counter each remove one of those guarantees in isolation so that the corresponding restriction can be studied without confronting all of them at once.
Chapter~2 will need the full combination, since its language will include constructs that can fail, interact, and loop.

\section{Fold only when the guard is known}
\label{sec:trap-guarded-folding}

The three extensions share one folding discipline.
A pure total subexpression may be evaluated during translation whenever its operands are known.
A subexpression that may trap, write, or diverge may be evaluated early only with an additional argument that the early evaluation preserves the program's behavior set.
Algorithm~\ref{alg:trap-guarded-fold} states the guarded rule for the division extension.
It traverses children first, as before, but consults the divisor before replacing a division.
The rule for addition and multiplication is unchanged because those operations remain pure and total in the extended language.
The rule would need further conditions if they could trap or write.

\begin{algorithm}
  \caption{Fold constants in the presence of a trapping division.}
  \label{alg:trap-guarded-fold}
  \begin{pseudocode}
    \Require A finite, non-sharing tree with $\operatorname{div}$ nodes
    \Ensure Same trap-or-value for every input $a\in V$
    \Function{GuardedFold}{$e$}
      \If{$e$ is a literal or the input node}
        \State \Return $e$
      \EndIf
      \State $p\gets\Call{GuardedFold}{\operatorname{left}(e)}$
      \State $q\gets\Call{GuardedFold}{\operatorname{right}(e)}$
      \If{$p=\operatorname{lit}(u)$ and $q=\operatorname{lit}(v)$}
        \If{$e$ is an addition}
          \State \Return $\operatorname{lit}((u+v)\bmod256)$
        \ElsIf{$e$ is a multiplication}
          \State \Return $\operatorname{lit}((u\times v)\bmod256)$
        \Else
          \If{$v=0$}
            \State \Return a division node with children $p,q$
          \Else
            \State \Return $\operatorname{lit}(u/v)$
          \EndIf
        \EndIf
      \EndIf
      \State \Return a node with $e$'s operator and children $p,q$
    \EndFunction
  \end{pseudocode}
\end{algorithm}

The correctness claim must also be extended.
For the pure language, Equation~\ref{eq:tiny-fold-correctness} equated two values.
For the extended language, the claim is that for every valid tree $e$ and every input $a$, the folded tree either traps exactly when the original traps or returns the same value.
The leaf cases are unchanged.
The addition and multiplication cases first propagate any child trap; when both children succeed, the earlier arithmetic argument applies.
For a division with folded children $p'$ and $q'$, the induction assumption gives trap-or-value agreement for each child.
If either child traps, both the retained division and the original division trap by Equation~\ref{eq:trap-div-semantics}.
If both succeed with $u,v$ and $v=0$, the algorithm retains the division, which traps as the original does.
If both succeed and $v\neq0$, the algorithm returns the literal $u/v$, which is exactly the retained division's value.
No case replaces a trap with a value or a value with a different value.

This proof covers the specified algorithm on valid trees.
It does not cover an implementation that forgets the $v=0$ test, evaluates the divisor with host division before testing it, or shares divisor nodes between two parents.
Each of those defects has a concrete shape in C.
Forgetting the test turns the always-trapping tree into a literal and is caught by the all-inputs trap check.
Dividing before testing executes undefined behavior on the host even when the source semantics calls for a defined trap.
Sharing a divisor node violates the non-sharing precondition on which both the linear-time argument and the in-place mutation reasoning depend.
The laboratory's storage rule therefore stays unchanged: the caller owns every node, disconnected nodes are not freed during the phase, and no candidate introduces sharing or cycles.
The new operation changes which replacements are legal, not how storage is managed.

The same separation of concerns from Section~\ref{sec:analysis-and-transformation} now pays for itself.
The fact \enquote{this divisor is the nonzero literal $8$} is an analysis result with a defined scope: a particular program point after folding.
The replacement of $3+5$ by $8$ uses that fact, and the retention of $(x+7)/0$ follows from the absence of the corresponding fact for a zero divisor.
If a later pass changes the divisor, the fact must be recomputed or discarded.
Continuing to use the old \enquote{nonzero} conclusion after the divisor has been rewritten would be the same mistake as continuing to use any invalidated analysis.
Stating the guard as part of the rule makes that dependency visible to the next pass.

Testing must also check the trap dimension explicitly.
The laboratory's division-pair test asserts trap-or-value agreement for all $65{,}536$ operand pairs, not merely value agreement where both sides succeed.
Its folding test asserts that a zero divisor keeps the division node rather than only asserting that some value is preserved.
Its composition is deliberately limited: the pure stack compiler is not asked to compile trapping trees, since its target has no checked instruction.
That limitation is stated as a contract boundary rather than left as a silent assumption.
A future target with a checked division instruction will need its own prefix theorem, verifier rule for stack depth on both success and failure paths, and composition proof before the two pipelines can be connected.
Safonov's treatment of compiler trustworthiness emphasizes that such boundaries are engineering obligations rather than incidental details~\cite{safonov-2010-trustworthy}.
Morgan's organization of representations around execution rules makes the same point from the implementation side: the IR must expose the facts each stage needs~\cite[Chapter~4]{morgan-1998-optimizing}.
Appel's sequence of intermediate languages shows how those exposures compose across a whole compiler~\cite{appel-1998-modern}.


\section{What we can now claim}
\label{sec:first-chapter-checkpoint}

We have defined a finite expression language, not a complete programming language.
Its evaluator follows four equations.
Its constant folder removes operations whose operands have become literals.
A structural proof establishes preservation for the specified algorithm, and executable tests check the implementation on exhaustive literal-operation inputs and selected tree shapes.
We also translate those trees to a four-operation stack machine.
A stack-prefix theorem connects target execution to source evaluation, while a separate verifier checks target structure and stack safety.
The combined pipeline folds a tree, emits code, verifies it, and executes it under explicit capacity limits.
An isolated faulty wrapper and a greedy reducer show how a structurally legal translation can still violate that pipeline's semantic promise.
A separate laboratory extends the tree with a trapping division, an ordered event log, and a fuel check.
It shows why a zero divisor must stay unfolded, why effect order belongs to the program's observations, and why divergence constrains code motion even when no wrong number appears.

We have also found the limits of that small result.
It does not justify deleting function effects, moving arbitrary operations across branches, or treating C's arithmetic as our expression arithmetic.
It does not establish that fewer tree nodes always means less elapsed time.
It does not remove the need for ownership rules or resource limits in the implementation.

Those limits are useful.
They tell us what additional information a more ambitious compiler will need.
Names and types need rules.
Branches need a graph.
Analyses need a meaning and a safe direction of approximation.
Transformations need contracts that remain valid when other passes change the program.
The rest of the project grows from those needs rather than from a list of clever rewrites.

\section{Exercises}
\label{sec:programs-meaning-exercises}

\begin{enumerate}
  \item \label{ex:tiny-wraparound}
  Evaluate $(x+7)\times(3+5)$ in the language of Equation~\ref{eq:tiny-expression-grammar} for $x=9$, $x=249$, and $x=255$.
  Show the reduction after each operation.
  Explain why a mathematical integer result outside $V$ is not a result of this language.

  \item \label{ex:tiny-nested-fold}
  Draw the expression tree for $(14+19)\times(6+5)$.
  Apply Algorithm~\ref{alg:tiny-fold} in postorder and record each replacement.
  Count the nodes reachable from the root before and after folding.

  \item \label{ex:tiny-zero-identity}
  Prove that $\operatorname{mul}(e,\operatorname{lit}(0))$ can be replaced by $\operatorname{lit}(0)$ in the finite, pure language defined here.
  State which properties of the language justify omitting evaluation of $e$.
  Explain why the same value identity does not justify deleting the call in Listing~\ref{lst:zero-effect}.

  \item \label{ex:tiny-identity-extension}
  Extend the folder with the rule $\operatorname{add}(e,\operatorname{lit}(0))\mapsto e$.
  Keep the caller-owned, non-sharing tree contract.
  Describe how an in-place implementation can replace the root node without returning a pointer to a temporary object.
  Test the rule for every input when $e$ is $x$ and when $e$ is $\operatorname{mul}(x,\operatorname{lit}(7))$.

  \item \label{ex:tiny-test-limits}
  A test compares original and folded results for all $256$ inputs on one tree and finds no difference.
  State exactly which preservation claim this establishes if the evaluator is correct.
  Give two claims about the folder that this test alone does not establish.

  \item \label{ex:tiny-shared-node}
  Suppose the no-sharing invariant is removed, and the same node is used as both children of a parent.
  Explain why the current linear-time argument in terms of distinct allocated nodes no longer follows from the tree traversal.
  Construct a family of acyclic shared graphs with few distinct nodes but many repeated recursive visits.
  Suggest one representation-aware way to avoid the repeated visits.
  \item \label{ex:tiny-stack-prefix}
  Execute the seven instructions in Table~\ref{tab:tiny-stack-trace} with input $9$ and initial stack $[23]$.
  Show every stack and identify the point of maximum depth.
  Explain why the final stack is correct for a fragment but not for a whole expression program started empty.

  \item \label{ex:tiny-stack-demand}
  Using Equation~\ref{eq:tiny-stack-demand}, compare the stack demands of these trees:
  \[
    \operatorname{add}(x,\operatorname{add}(x,x)),
    \qquad
    \operatorname{add}(\operatorname{add}(x,x),x).
  \]
  Check the arithmetic equivalence in the modulo-256 language, and state why this observation alone would not justify reassociation in an arbitrary effectful or floating-point language.

  \item \label{ex:tiny-stack-valid-wrong}
  Give a target program that passes the verifier but is not a correct translation of $x$ for every input.
  Give an input that exposes the mismatch.
  Identify which obligation the verifier establishes and which it does not.
  \item \label{ex:tiny-mutant-hidden}
  Fold the reduced counterexample $0\times1$ before compiling it with the faulty wrapper.
  Explain why the mismatch disappears and why this is not a repair of the wrapper.

  \item \label{ex:tiny-mutant-agreement}
  Find all pairs $u,v\in V$ for which translating one multiplication as addition accidentally gives the correct result.
  Derive a congruence for those pairs, and determine the number of pairs.
  Explain why testing only $1\times1$ would still detect the mutation, whereas testing only $0\times0$ would not.

  \item \label{ex:trap-div-zero}
  Evaluate $(x+7)/(3+5)$ in the extended language of Equation~\ref{eq:trap-div-semantics} for $x=9$ and $x=250$.
  Show each reduction, including the folded divisor $8$ and the truncated quotients $2$ and $0$.
  Then consider $(x+7)/(249+7)$.
  Show that the divisor folds to $0$, explain why the whole expression traps for every input, and state why replacing it with any literal in $V$ would be wrong.

  \item \label{ex:trap-effect-swap}
  Two calls return $10$ and $20$ while logging tags $1$ and $2$ in call order.
  Show that left-to-right evaluation logs $[1,2]$ and right-to-left evaluation logs $[2,1]$, with sum $30$ in both cases.
  Explain why swapping the calls preserves the value but changes the trace, and which contract decides whether the swap is allowed.

  \item \label{ex:trap-hoist-branch}
  Consider a branch that divides $10$ by $b$ only when $b\neq0$, and otherwise returns $0$.
  Give concrete values with $b=0$ where the guarded program returns normally.
  Show that hoisting the division above the test as in Figure~\ref{fig:trap-hoist} introduces a trap on that path.
  State the resulting refinement failure using Equation~\ref{eq:behavior-refinement}.
\end{enumerate}

\section{Further reading}
\label{sec:programs-meaning-reading}

Appel's \emph{Modern Compiler Implementation in ML} connects front-end and back-end algorithms through an implementable compiler organization~\cite{appel-1998-modern}.
Read it for the way interfaces and data structures give the individual algorithms a place in a working system.

Morgan's \emph{Building an Optimizing Compiler} develops the representations and analyses needed to connect optimization to code generation~\cite{morgan-1998-optimizing}.
Its discussion of execution rules is a useful reminder that syntax alone does not define the transformations a language permits.

Safonov's \emph{Trustworthy Compilers} broadens the discussion to the engineering obligations of dependable translation~\cite{safonov-2010-trustworthy}.
Use it to compare different kinds of assurance rather than treating successful compilation as evidence that the compiler is correct.

The WG14 C11 committee draft is a specification, not a compiler tutorial~\cite{wg14-2011-c11}.
For the examples here, consult its environmental limits, unsigned arithmetic rules, expression rules, and abstract-machine requirements.
Checking those clauses against the code is more useful than assuming that a familiar C type behaves like a mathematical number.
